\SGAchapter{XV}{群预概形的子环面补论；在光滑群中的应用}{Compléments sur les sous-tores d'un préschéma en groupes. Application aux groupes lisses}
\SourcePage{255}
\begin{center}M.~Raynaud\end{center}
\footnote{原文版本记号：2008 年 12 月 1 日 xy 版。}
\footnote{本讲及后两讲（第 XVI、XVII 讲）并不对应于讨论班上的口头报告。它们在作出若干补充的同时，展开了 A.~Grothendieck 于 1964 年为本讨论班撰写的简略手稿笔记的实质内容。}
\SGAsection{0}{引言}
本讲补充第 XI、XII 讲，并部分重述其中的内容；不必预先了解第 XIII、XIV 讲。我们继续第 XII 讲开始的工作，研究不一定在 $S$ 上\Term{affine}{仿射}、也不一定在 $S$ 上\Term{separated}{分离}的 $S$-\Term{group-prescheme}{群预概形}。

第 1、2、3、4 节研究\Term{group-prescheme}{群预概形}的\Term{subtorus}{子环面}。我们将得到\Term{infinitesimal}{无穷小}\Term{lifting}{提升}\Term{theorem}{定理}（第 2 节）和整体\Term{lifting}{提升}\Term{theorem}{定理}（第 4 节），其中\Term{finite-order-point}{有限阶点}（第 1 节）起着关键\Term{action}{作用}。

第 5、6、7 节独立于前面各节。考察\Term{infinitesimal-neighborhood}{无穷小邻域}，可以证明那些等于自身\Term{connected-normalizer}{连通正规化子}的\Term{smooth}{光滑}\Term{subgroup}{子群}所成\Term{functor}{函子}的\Term{representability}{可表性}（第 5 节）。第 6、7 节特别关注 \Term{cartan-subgroup}{Cartan 子群}。

最后，第 8 节给出\Term{smooth}{光滑}群的\Term{subtorus}{子环面}\Term{functor}{函子}以及\Term{maximal-torus}{极大环面}\Term{functor}{函子}\Term{representability}{可表}的一个充要条件。
\SGAsection{1}{有限子群的提升}
\textbf{1. \Term{finite}{有限}、\Term{smooth}{光滑}且\Term{central}{中心的}\Term{multiplicative-type}{乘法型}\Term{subgroup}{子群}}
\begin{SGAresult}{命题}{1.1}
设 $S$ 是\Term{affine}{仿射}\Term{scheme}{概形}，$S_0$ 是由平方为零的\Term{ideal}{理想}定义的 $S$ 的\Term{closed-subscheme}{闭子概形}，$G$ 是 $S$-\Term{group-prescheme}{群预概形}，$H_0$ 是 $G_0=G\times_S S_0$ 的一个\Term{group-scheme}{子群概形}，它在 $S_0$ 上\Term{smooth}{光滑}、\Term{finite-type}{有限型}且为\Term{multiplicative-type}{乘法型}。则存在 $G$ 的一个\Term{multiplicative-type}{乘法型}\Term{group-scheme}{子群概形}\Term{lifting}{提升} $H_0$，
\SourcePage{256}
当且仅当存在 $G$ 的一个在 $S$ 上\Term{flat}{平坦}的\Term{subscheme}{子概形} $H'$ \Term{lifting}{提升} $H_0$。
\end{SGAresult}
\begin{proof}
必要性显然，来证充分性。\Term{multiplicative-type}{乘法型}群 $H_0$ 是\Term{quasi-isotrivial}{拟等平凡}的（第 X 讲 4.5）；根据第 X 讲 2.1，存在一个\Term{multiplicative-type}{乘法型} $S$-群 $H$ 及一个 $S_0$-群\Term{isomorphism}{同构}
\[u_0:H\times_S S_0\xrightarrow{\sim}H_0.\]
由于 $H'$ 在 $S$ 上\Term{flat}{平坦}，而 $H'\times_S S_0$ 在 $S_0$ 上\Term{finite-presentation}{有限呈示}（第 IX 讲 2.1 b)），$H'$ 在 $S$ 上\Term{finite-presentation}{有限呈示}；此外，其\Term{fiber}{纤维}\Term{smooth}{光滑}，所以 $H'$ 在 $S$ 上\Term{smooth}{光滑}（EGA IV 17.5.1）。由于\Term{scheme}{概形} $S$ \Term{affine}{仿射}，$u_0$ \Term{liftable}{可提升}为一个\Term{prescheme}{预概形}的 $S$-\Term{morphism}{态射}
\[u:H\longrightarrow H'.\]
由第 III 讲 2.1 及第 IX 讲 3.1，\Term{composition-morphism}{复合态射}
\[v_0:H\times_{S_0}S\xrightarrow{\sim}H_0\longrightarrow G_0\]
也\Term{liftable}{可提升}为一个 $S$-\Term{group-morphism}{群态射}
\[v:H\longrightarrow G.\]
由于 $v_0$ 是\Term{immersion}{浸入}，$v$ 也是\Term{immersion}{浸入}。因此，$H$ 在 $v$ 下的像是 $G$ 的一个\Term{multiplicative-type}{乘法型}\Term{group-scheme}{子群概形}，并且\Term{lifting}{提升}了 $H_0$。
\end{proof}
\begin{SGAresult}{命题}{1.2}
设 $S$ 是\Term{prescheme}{预概形}，$S_0$ 是由一个\Term{locally-nilpotent}{局部幂零}\Term{ideal-sheaf}{理想层}定义的 $S$ 的\Term{subprescheme}{子预概形}，$G$ 是在 $S$ 上\Term{flat}{平坦}且\Term{finite-presentation}{有限呈示}的 $S$-\Term{group-prescheme}{群预概形}，$H_0$ 是 $G_0=G\times_S S_0$ 的一个\Term{group-scheme}{子群概形}，它在 $S_0$ 上\Term{smooth}{光滑}且\Term{finite}{有限}，并且是\Term{multiplicative-type}{乘法型}的\Term{central}{中心}\Term{subgroup}{子群}。则存在唯一的 $G$ 的\Term{multiplicative-type}{乘法型}\Term{group-scheme}{子群概形} $H$ \Term{lifting}{提升} $H_0$。此外，$H$ 是\Term{central}{中心}\Term{subgroup}{子群}。（见第 XVII 讲附录 III，1。）
\end{SGAresult}
\begin{SGAresult}{命题}{1.2 bis}
设 $S,G,S_0,G_0$ 如上，$H$ 是一个\Term{multiplicative-type}{乘法型}、在 $S$ 上\Term{smooth}{光滑}且\Term{finite}{有限}的 $S$-\Term{group-scheme}{群概形}，$u_0:H_0=H\times_S S_0\to G_0$ 是一个\Term{central}{中心}\Term{homomorphism}{同态}。则 $u$ 唯一地\Term{lifting}{提升}为一个\Term{homomorphism}{同态} $u:H\to G$。此外，$u$ 是\Term{central}{中心}\Term{homomorphism}{同态}。
\end{SGAresult}
考察 $u_0$ 的图，容易从 1.2 \Term{pushout}{推出} 1.2 bis 中\Term{lifting}{提升} $u$ 的存在性。由第 IX 讲 3.4 和 5.1，\Term{lifting}{提升} $u$ 唯一且为\Term{central}{中心}\Term{homomorphism}{同态}。

\textbf{1.2 的证明。}由第 IX 讲 5.6 b) 及 3.4 bis，$H$ 唯一，且 $H$ 为\Term{central}{中心}\Term{subgroup}{子群}。利用唯一性，为证明 $H$ 存在，可以假定 $S$ \Term{affine}{仿射}，且 $S_0$ 由一个平方为零的\Term{ideal}{理想}定义；根据 1.1，只需找到一个在 $S$ 上\Term{flat}{平坦}的 $G$ 的\Term{subscheme}{子概形}\Term{lifting}{提升} $H_0$。

由于 $H_0$ 在 $S_0$ 上\Term{smooth}{光滑}且\Term{finite}{有限}，缩小 $S$ 后可以假定存在一个在 $S$ 上\Term{invertible}{可逆}的整数 $n>0$，使 $H_0={}_nH_0$。考察 $G$ 中的 $n$ 次幂\Term{morphism}{态射}
\[u:G\longrightarrow G,\qquad x\longmapsto x^n.\]
\SourcePage{257}
仍用 ${}_nG$ 表示“$u$ 的核”，即 $G$ 的\Term{unit-section}{单位截面}在 $u$ 下的\Term{inverse-image}{逆像}（注意，$u$ 一般不是\Term{group-morphism}{群态射}）。暂且承认下面的\Term{lemma}{引理}。
\begin{SGAresult}{\Term{lemma}{引理}}{1.3}
设 $k$ 是域，$G$ 是 $k$ 上\Term{locally-finite-type}{局部有限型}的\Term{group-scheme}{群概形}，$n$ 是与 $k$ 的\Term{characteristic}{特征}互素的整数，$u$ 是 $G$ 中的 $n$ 次幂\Term{morphism}{态射}。则 $u$ 在 $G$ 的\Term{center}{中心}内的每一点 $x$ 处\Term{etale}{平展}。
\end{SGAresult}
由于 $G$ 在 $S$ 上\Term{flat}{平坦}且\Term{finite-presentation}{有限呈示}，由上面的\Term{lemma}{引理}及 EGA IV 17.8.2，若 $G$ 的一点 $x$ \Term{projection}{投影}到 $S$ 的点 $s$，且属于 $G_s$ 的\Term{center}{中心}，则\Term{morphism}{态射} $u$ 在 $x$ 处\Term{etale}{平展}。若 $x$ 还属于 ${}_nG$，则 ${}_nG$ 在 $x$ 处对 $S$ \Term{etale}{平展}。根据假设，群 $H_0$ 是\Term{central}{中心}\Term{subgroup}{子群}且包含在 ${}_nG_0$ 中，因此实际上包含在 ${}_nG$ 的对 $S$ \Term{etale}{平展}的最大\Term{open-subset}{开集} $V$ 中。由于 $H_0$ 和 $V\times_S S_0=V_0$ 都在 $S_0$ 上\Term{etale}{平展}，$H_0$ 是 $V_0$ 的一个\Term{open-subset}{开集}（EGA IV 17.8.7 及 17.9.1）。于是，$V$ 中与 $H_0$ 具有相同底\Term{space}{空间}的\Term{open-subprescheme}{开子预概形}就是 $G$ 的一个在 $S$ 上\Term{flat}{平坦}的\Term{subscheme}{子概形}，它\Term{lifting}{提升}了 $H_0$。

还须证明\Term{lemma}{引理} 1.3。注意，$G$ 上 $u$ 为\Term{etale}{平展}\Term{morphism}{态射}的最大\Term{open-subset}{开集}在\Term{base-field-extension}{基域扩张}下保持\Term{invariant}{不变}（EGA IV 17.8.2）；这使我们可以归结到 $x$ 在 $k$ 上有理的情形。以 $t_x$ 表示由 $x$ 定义的\Term{translation}{平移}，它是\Term{scheme}{概形} $G$ 的一个 $k$-\Term{automorphism}{自同构}。由于 $x$ 位于 $G$ 的\Term{center}{中心}，有\Term{commutative-diagram}{交换图}
\[
\begin{tikzcd}
G\arrow[r,"u"]\arrow[d,"t_x"']&G\arrow[d,"t_x^n"]\\
G\arrow[r,"u"']&G.
\end{tikzcd}
\]
因此只需证明 $u$ 在原点\Term{etale}{平展}，而这一点已在第 VII$_A$ 讲 8.4 中证明。

\textbf{2. \Term{finite-group}{有限群}的整体\Term{lifting}{提升}}
\begin{SGAresult}{\Term{lemma}{引理}}{1.4}
设 $A$ 是\Term{local-ring}{局部环}，对于由其\Term{maximal-ideal}{极大理想} $\mathfrak m$ 定义的\Term{topology}{拓扑}\Term{separated}{分离}且\Term{complete}{完备}，令 $S=\operatorname{Spec}(A)$，$S_n=\operatorname{Spec}(A/\mathfrak m^n)$。则对任意\Term{prescheme}{预概形} $X$（相应地，任意 $S$-\Term{prescheme}{预概形} $X$），\Term{canonical}{典范}映射
\[\operatorname{Hom}(S,X)\longrightarrow\varprojlim_n\operatorname{Hom}(S_n,X)\tag{x}\]
（相应地，
\[\Gamma(X/S)\longrightarrow\varprojlim_n\Gamma(X_n/S_n),\qquad X_n=X\times_S S_n\tag{xx}\]
）是\Term{bijection}{双射}。
\end{SGAresult}
(xx) 容易由 (x) \Term{pushout}{推出}。下面证明 (x)。设 $u_n:S_n\to X$（$n\in\mathbb N$）是一族相容的\Term{morphism}{态射}。$S_n$ 的\Term{closed-point}{闭点}的像 $y$ 因而与 $n$ 无关，且 $u_n$ 经由 $\operatorname{Spec}\mathcal O_y$ \Term{resolution}{分解}。
\SourcePage{258}
对\Term{morphism}{态射} $u_n$ 取\Term{inverse-limit}{投射极限}，便定义一个\Term{ring-homomorphism}{环同态}
\[\widetilde u:\mathcal O_y\longrightarrow\varprojlim_n(A/\mathfrak m^n)=A.\]
这说明 (x) 为\Term{surjectivity}{满射}；只要 $A$ 对 $\mathfrak m$-\Term{adic-topology}{进拓扑}\Term{separated}{分离}，它就是\Term{injectivity}{单射}。
\begin{SGAresult}{推论}{1.5}
设 $A$ 是\Term{complete}{完备} noether \Term{local-ring}{局部环}，$\mathfrak m$ 是其\Term{maximal-ideal}{极大理想}，$S=\operatorname{Spec}(A)$，$S_n=\operatorname{Spec}(A/\mathfrak m^n)$，$X$ 是 $S$ 上的\Term{finite}{有限}\Term{scheme}{概形}，$Y$ 是 $S$-\Term{prescheme}{预概形}。则\Term{canonical}{典范}映射
\[\operatorname{Hom}_S(X,Y)\longrightarrow\varprojlim_n\operatorname{Hom}_{S_n}(X_n,Y_n)\]
是\Term{bijection}{双射}，其中 $X_n=X\times_S S_n$，同样 $Y_n=\cdots$。
\end{SGAresult}
事实上，由 EGA II 6.2.5，$X$ 是\Term{finite}{有限}多个\Term{local-scheme}{局部}的、在 $S$ 上\Term{finite}{有限}的 $S$-\Term{scheme}{概形}的和。这使我们归结到 $X$ 本身是一个\Term{complete}{完备} noether \Term{local-ring}{局部环}的谱的情形。但 $\operatorname{Hom}_S(X,Y)=\Gamma(Z/X)$，其中 $Z$ 是 $X$-\Term{prescheme}{预概形} $Y\times_S X$，于是应用上面的命题即可。
\begin{SGAresult}{命题}{1.6}
设 $A,S,S_n$ 如上，$G,M$ 是两个 $S$-\Term{group-prescheme}{群预概形}，其中 $M$ 在 $S$ 上\Term{finite}{有限}。则：

a) \Term{canonical}{典范}映射
\[\operatorname{Hom}_{S\text{-gr}}(M,G)\longrightarrow\varprojlim_n\operatorname{Hom}_{S_n\text{-gr}}(M_n,G_n)\]
为\Term{bijection}{双射}。

b) 若 $M$ 为\Term{multiplicative-type}{乘法型}且 $G$ 在 $S$ 上\Term{smooth}{光滑}，则\Term{canonical}{典范}映射
\[\varphi:\operatorname{Hom}_{S\text{-gr}}(M,G)\longrightarrow\operatorname{Hom}_{S_0\text{-gr}}(M_0,G_0)\]
为\Term{surjectivity}{满射}。此外，若 $\varphi(u)=\varphi(u')=u_0$，则 $u,u'$ 由 $G(S)$ 中一个\Term{reduced}{约化}为 $G(S_0)$ 的单位元素的元素\Term{conjugacy}{共轭}。

c) 若 $M$ 是\Term{multiplicative-type}{乘法型}且在 $S$ 上\Term{smooth}{光滑}，$G$ 在 $S$ 上\Term{flat}{平坦}且\Term{finite-type}{有限型}，$u_0:M_0\to G_0$ 是一个\Term{central}{中心}\Term{homomorphism}{同态}，则 $u_0$ 唯一地\Term{lifting}{提升}为一个\Term{homomorphism}{同态}
\[u:M\longrightarrow G.\]
此外，若 $G$ 的\Term{fiber}{纤维}\Term{connected}{连通}，则 $u$ 为\Term{central}{中心}\Term{homomorphism}{同态}。
\end{SGAresult}
\begin{proof}
a) 由 1.5、$M\times_S M$ 在 $S$ 上\Term{finite}{有限}，以及\Term{group-morphism}{群态射}是使下面这个熟知的图交换的\Term{morphism}{态射}这一刻画得到：
\[
\begin{tikzcd}
M\times_S M\arrow[r,"u\times u"]\arrow[d]&G\times_S G\arrow[d]\\
M\arrow[r,"u"']&G.
\end{tikzcd}
\]
\SourcePage{259}
b) 根据第 IX 讲 3.6，可以构造一族相容的\Term{homomorphism}{同态} $u_n:M_n\to G_n$，它们\Term{lifting}{提升}给定的\Term{homomorphism}{同态} $u_0:M_0\to G_0$。结合 a)，便得到 b) 的第一个断言。

现若 $u,u'$ 是 $u_0$ 的两个\Term{lifting}{提升}，则 $u_n,u'_n$ 由 $G(S_n)$ 中一个\Term{lifting}{提升} $G(S_0)$ 的单位元素的元素 $g_n$ \Term{conjugacy}{共轭}（第 IX 讲 3.6）；同处的结果还说明，可以相容地选择 $g_n$，于是由 1.4，它们来自 $G(S)$ 的一个\Term{section}{截面} $g$。\Term{morphism}{态射} $u$ 和 $\operatorname{int}(g)u'$ 模每个 $\mathfrak m^n$ 都相等，因而相等（1.5）。

c) $u$ 的存在性和唯一性由 a) 与 1.2 bis 得到。若 $G$ 的\Term{fiber}{纤维}\Term{connected}{连通}，则由第 IX 讲 5.6 a)，$u$ 是\Term{central}{中心}\Term{homomorphism}{同态}。
\end{proof}
\SGAsection{2}{子环面的无穷小提升}
\textbf{1. \Term{theorem}{定理}的陈述}

我们将给出一个\Term{group-prescheme}{群预概形}的\Term{subtorus}{子环面}的\Term{infinitesimal}{无穷小}\Term{lifting}{提升}\Term{theorem}{定理}。它不要求\Term{smooth}{光滑}性假设（与第 IX 讲 3.6 bis 不同），同时回答一个很自然的问题\footnote{编者注：用\Term{finite}{有限}\Term{subgroup}{子群}逼近一个\Term{torus}{环面}的想法，已出现在 Grothendieck 对\Term{torus}{环面}的\Term{centralizer}{中心化子}\Term{connected}{连通}性的证明中，见 BIBLE 第 4.6 节。}：只要能够\Term{lifting}{提升}\Term{subtorus}{子环面} $T_0$ 的“足够多”的\Term{finite-order-point}{有限阶点}，是否就能保证\Term{lifting}{提升} $T_0$（当然是在\Term{infinitesimal}{无穷小}意义下）？
\begin{SGAresult}{\Term{theorem}{定理}}{2.1}
设 $S$ 是\Term{affine}{仿射} noether \Term{scheme}{概形}，$S_0$ 是由平方为零的\Term{ideal}{理想} $J$ 定义的 $S$ 的\Term{closed-subscheme}{闭子概形}，$G$ 是\Term{finite-type}{有限型}的 $S$-\Term{group-prescheme}{群预概形}，$G_0=G\times_S S_0$，$T_0$ 是 $G_0$ 的\Term{subtorus}{子环面}，$q$ 是在 $S$ 上\Term{invertible}{可逆}的整数，$q>0$。假定对每个等于 $q$ 的幂的整数 $n$，存在 $G$ 的一个\Term{subscheme}{子概形} $M_n$，它在 $S$ 上\Term{flat}{平坦}且满足 $M_n\times_S S_0={}_nT_0$。则存在 $G$ 的一个\Term{subtorus}{子环面} $T$，使 $T\times_S S_0=T_0$。
\end{SGAresult}
我们将通过下面两个推论使用\Term{theorem}{定理} 2.1。
\begin{SGAresult}{推论}{2.2}
设 $S$ 是\Term{locally-noetherian}{局部 noether} \Term{prescheme}{预概形}，$S_0$ 是由\Term{locally-nilpotent}{局部幂零}\Term{ideal-sheaf}{理想层}定义的 $S$ 的\Term{closed-subprescheme}{闭子预概形}，$G$ 是\Term{finite-type}{有限型}的 $S$-\Term{group-prescheme}{群预概形}，$T_0$ 是 $G_0=G\times_S S_0$ 的\Term{subtorus}{子环面}，$q$ 是在 $S$ 上\Term{invertible}{可逆}的整数，$q>0$。最后，当整数 $n$ 遍历 $q$ 的幂时，设 $(M_n)$ 是一族相容的、\Term{lifting}{提升} ${}_nT_0$ 的 $G$ 的\Term{multiplicative-type}{乘法型} $S$-\Term{group-scheme}{子群概形}。（注意，称\Term{multiplicative-type}{乘法型}\Term{subgroup}{子群}族 $(M_n)$ 相容，是指当整数 $m$ 整除 $n$ 时 $M_m={}_m(M_n)$。）则恰有一个 $G$ 的\Term{subtorus}{子环面} $T$，使 $T\times_S S_0=T_0$，且对每个 $n$，有 ${}_nT=M_n$。
\end{SGAresult}
\begin{SGAresult}{推论}{2.3}
设 $G$ 是在 $S$ 上\Term{flat}{平坦}且\Term{finite-presentation}{有限呈示}的 $S$-\Term{group-prescheme}{群预概形}，$S_0$ 是由一个\Term{finite-type}{有限型}且\Term{locally-nilpotent}{局部幂零}的\Term{ideal-sheaf}{理想层}定义的 $S$ 的\Term{closed-subprescheme}{闭子预概形}，$T_0$ 是 $G_0=G\times_S S_0$ 的\Term{central}{中心}\Term{torus}{环面}。则恰有一个 $G$ 的\Term{subtorus}{子环面} $T$ \Term{lifting}{提升} $T_0$。此外，$T$ 是\Term{central}{中心}\Term{subgroup}{子群}。
\end{SGAresult}
\SourcePage{260}
\textbf{注记 2.4。}我们将 2.1、2.2、2.3 的下述类似表述留给读者：不再\Term{lifting}{提升} $G_0$ 的一个\Term{subtorus}{子环面}，而是给定 $S$ 上的一个\Term{torus}{环面} $T$，要求\Term{lifting}{提升}一个\Term{morphism}{态射}
\[u_0:T_0\longrightarrow G_0\]
（考察 $u_0$ 的图，便归结到前面的情形）。

现在说明如何由 2.1 \Term{pushout}{推出}推论 2.2 和 2.3。

\textbf{推论 2.2 的证明。}
$T$ 的唯一性由第 IX 讲 4.8 b) 和 4.10 得到。因此，为证明 $T$ 存在，可以归结到 $S$ \Term{affine}{仿射}、从而 noether，且 $S_0$ 由平方为零的\Term{ideal}{理想}定义的情形。
\begin{SGAresult}{\Term{lemma}{引理}}{2.5}
设 $G$ 是\Term{finite-presentation}{有限呈示}的 $S$-\Term{group-prescheme}{群预概形}，$H$ 是 $G$ 的一个\Term{multiplicative-type}{乘法型}、在 $S$ 上\Term{smooth}{光滑}的\Term{group-scheme}{子群概形}。则 $\operatorname{Centr}_G(H)$ 可由 $G$ 的一个\Term{finite-presentation}{有限呈示}的\Term{group-prescheme}{子群预概形}表示。
\end{SGAresult}
当 $G$ 在 $S$ 上\Term{separated}{分离}时，本\Term{lemma}{引理}由第 VIII 讲 6.5 e) 得到，而无须假设 $H$ \Term{smooth}{光滑}。在这里，注意所要证明的断言对于 \Term{fpqc-topology}{fpqc 拓扑}是\Term{local-scheme}{局部}的，因此可假定 $H$ \Term{diagonalizable}{可对角化}，从而具有 $\mathrm D_S(M)$ 的\Term{form}{形式}。还可假定 $S$ \Term{affine}{仿射}，再由 EGA IV 8 假定 $S$ noether。由于 $H$ 在 $S$ 上\Term{smooth}{光滑}且\Term{finite-type}{有限型}，$M$ 的\Term{torsion-subgroup}{挠子群}的阶 $q$ 在 $S$ 上\Term{invertible}{可逆}（第 VIII 讲 2.1 e)）。于是立即可知（参见第 IX 讲 4.10），当 $n$ 遍历 $q$ 的幂时，\Term{subgroup}{子群} ${}_nH$ 在 $H$ 中\Term{schematically-dense}{概形稠密}（第 IX 讲 4.1）。但 ${}_nH$ 是 $S$ 的一个\Term{completely-split-covering}{完全分解覆盖}（即\Term{isomorphism}{同构}于\Term{finite}{有限}个 $S$ 的副本的\Term{direct-sum}{直和}），所以 $\operatorname{Centr}_G({}_nH)=Z_n$ 可由 ${}_nH$ 在 $S$ 上的各个\Term{section}{截面}在 $G$ 中的\Term{centralizer}{中心化子}之交表示。于是只需应用下面的\Term{lemma}{引理}。
\begin{SGAresult}{\Term{lemma}{引理}}{2.5 bis}
设 $S$ 是 noether \Term{prescheme}{预概形}，$G$ 是\Term{finite-type}{有限型}的 $S$-\Term{group-prescheme}{群预概形}，$H$ 是 $G$ 的一个\Term{multiplicative-type}{乘法型}\Term{subgroup}{子群}，$H_i$ 是 $G$ 的一族递增\Term{filtered-category}{滤过的}\Term{multiplicative-type}{乘法型}\Term{subgroup}{子群}，并假定 $Z_i=\operatorname{Centr}_G(H_i)$ 可由 $G$ 的\Term{group-prescheme}{子群预概形}表示。则 $Z_i$ 族最终\Term{stable}{稳定}。若 $H_i$ 还在 $H$ 中\Term{schematically-dense}{概形稠密}，则对充分大的 $i$，有 $Z_i=\operatorname{Centr}_G(H)$。
\end{SGAresult}
要证明 $Z_i$ 族最终\Term{stable}{稳定}，只需证明其\Term{underlying-set}{底集}合族 $\operatorname{ens}(Z_i)$ 最终\Term{stable}{稳定}。事实上，其\Term{stable}{稳定}值是 $G$ 的一个\Term{open-subset}{开集} $U$ 中的\Term{closed-subset}{闭子集}，故以 $U$ 代替 $G$，便归结到研究一个 noether \Term{prescheme}{预概形}的一族递减\Term{filtered-category}{滤过的}\Term{closed-subprescheme}{闭子预概形}。一个简单的\Term{constructibility}{可构造性}论证将我们归结到 $S$ \Term{integral}{整}的情形。此时须证明 $\operatorname{ens}(Z_i)$ 族在 $S$ 的某个\Term{nonempty}{非空}\Term{open-subset}{开集}上最终\Term{stable}{稳定}。但 $G$ 的\Term{generic-fiber}{泛纤维}\Term{separated}{分离}（第 VI$_A$ 讲 0.3），因此缩小 $S$ 后，可以假定 $G$ 在 $S$ 上\Term{separated}{分离}（EGA IV 8）。此时 $Z_i$ 在 $G$ 中为闭（第 VIII 讲 6.5 e)）。
\SourcePage{261}
为证明\Term{lemma}{引理}的最后一个断言，记 $Z$ 为 $Z_i$ 族的\Term{stable}{稳定}值。显然，$\operatorname{Centr}_G(H)$ 是 $Z$ 的子\Term{functor}{函子}；下面证明 $Z$ \Term{centralize}{中心化} $H$。设 $E$ 是 $H\times_S Z$ 中作为下述两个\Term{morphism}{态射}的\Term{equalizer}{等化子}的\Term{subprescheme}{子预概形}：
\[\begin{aligned}
H\times_S Z&\rightrightarrows G,\\
(h,c)&\longmapsto c,\\
(h,c)&\longmapsto hch^{-1}.
\end{aligned}\]
\Term{prescheme}{预概形} $E$ 对每个 $i$ 都包含 $H_i\times_S Z$。另一方面，$H_i$ 在 $S$ 上\Term{flat}{平坦}，所以对 $S$ 的每一点 $s$，$(H_i)_s$ 在 $H_s$ 中\Term{schematically-dense}{概形稠密}，而 $H_i\times_S Z$ 在 $H\times_S Z$ 中\Term{schematically-dense}{概形稠密}（EGA IV 11.10.9）。由于 $G_s$ \Term{separated}{分离}，$E_s$ 在 $(H\times_s Z)_s$ 中为闭，因而与之相等。于是 $E$ 在 $H\times_S Z$ 中为闭，从而等于 $H\times_S Z$。这就是说，$Z$ \Term{centralize}{中心化} $H$。

回到 2.2 的证明。由 2.5，$Z_n=\operatorname{Centr}_G(M_n)$ \Term{representability}{可表}，由 2.5 bis，\Term{subprescheme}{子预概形}的递减族 $Z_n$ 最终\Term{stable}{稳定}，设其\Term{stable}{稳定}值为 $Z$。群 $Z$ 包含 $T_0$ 及各个 $M_n$。因此以 $Z$ 代替 $G$，便可假定各 $M_n$ 为\Term{central}{中心}\Term{subgroup}{子群}。

现在满足\Term{theorem}{定理} 2.1 的适用条件，故存在 $G$ 的一个\Term{subtorus}{子环面} $T$ \Term{lifting}{提升} $T_0$。群 ${}_nT$ 和 $M_n$ 是 ${}_nT_0$ 的两个\Term{lifting}{提升}，因而\Term{conjugacy}{共轭}（第 IX 讲 3.2 bis）；又因为 $M_n$ 为\Term{central}{中心}\Term{subgroup}{子群}，二者相等。\Term{torus}{环面} $T$ 即为所求。

\textbf{推论 2.3 的证明。}
$T$ 的唯一性由第 IX 讲 5.1 bis 得到，而 $T$ 为\Term{central}{中心}\Term{subgroup}{子群}则由第 IX 讲 5.6 b) 得到。利用这一点，按通常方法可以归结到 $S$ \Term{affine}{仿射}（因而 $S_0$ 由\Term{finite-type}{有限型}\Term{nilpotent-ideal}{幂零理想}定义）的情形，进而归结到 $S$ noether 的情形。必要时缩小 $S$，可以假定存在一个在 $S$ 上\Term{invertible}{可逆}的整数 $q$。于是推论 2.3 由 2.2 和 1.2 得到。

\textbf{注记 2.6。}容易证明，将\Term{torus}{环面} $T_0$ 换成 $G_0$ 的一个\Term{smooth}{光滑}\Term{central}{中心}\Term{multiplicative-type}{乘法型}\Term{subgroup}{子群}后，推论 2.3 仍成立。

\textbf{2. 2.1 的证明}

\textbf{a) 归结到 $T_0=\mathbb G_{m,S_0}$ 的情形。}
由 1.1，可以假定 $M_n$ 是\Term{multiplicative-type}{乘法型}\Term{subgroup}{子群}。利用第 IX 讲 3.2 bis，可以假定 $M_n$ 族相容（2.2）。$M_n$ 在 $G$ 中的\Term{centralizer}{中心化子} $Z_n$ \Term{representability}{可表}（2.5），且 $Z_n$ 的\Term{filtered-family}{滤过族}最终\Term{stable}{稳定}（2.5 bis）。因此，以充分大的 $n$ 对应的 $Z_n$ 代替 $G$，可以假定 $T_0$ 及各 $M_n$ 为\Term{central}{中心}\Term{subgroup}{子群}。此时 $T_0$ 的\Term{lifting}{提升}的唯一性得到保证（第 IX 讲 5.1 bis）。

像 1.1 的证明一样，可以假定存在一个 $S$-\Term{torus}{环面} $T$ 及一个 $S_0$-\Term{isomorphism}{同构}
\[u_0:T\times_S S_0\xrightarrow{\sim}T_0,\]
\SourcePage{262}
而\Term{lifting}{提升} $u_0$ 与\Term{lifting}{提升} $T_0$ \Term{equivalence}{等价}。由唯一性，只需在作\Term{faithfully-flat}{忠实平坦}、\Term{affine}{仿射}且\Term{finite-type}{有限型}的\Term{group-extension}{扩张} $S'\to S$ 后证明 $u_0$ 的\Term{lifting}{提升}存在（fpqc \Term{descent}{下降}），因而可以假定 $T=\mathbb G_{m,S}^r$（第 X 讲 4.5）。若 $u_0$ 在每个因子 $\mathbb G$ 上的\Term{restriction}{限制}均\Term{lifting}{提升}为一个 $S$-\Term{morphism}{态射}（它必为\Term{central}{中心}\Term{morphism}{态射}），便立即得到 $u_0$ 的一个\Term{lifting}{提升}。总之，可以假定 $T_0=\mathbb G_{m,S_0}$。

\textbf{b) 定义 $T_0$ 存在\Term{lifting}{提升}的\Term{obstruction}{障碍}。}
为证明 2.1，由 1.1，只需找到 $G$ 的一个在 $S$ 上\Term{flat}{平坦}的\Term{subscheme}{子概形}\Term{lifting}{提升} $T_0$。我们将看到，这种\Term{lifting}{提升}的存在\Term{obstruction}{障碍}可定义为某个模\Term{sheaf}{层}的 $\operatorname{Ext}^1(\ ,\ )$ 中的元素。

设 $U$ 是 $G$ 的一个\Term{open-subset}{开集}，使 $T_0$ 在 $U$ 中为闭；仍以 $U$（相应地 $U_0$）表示 $G$（相应地 $G_0$）中以 $U$ 为底\Term{space}{空间}的\Term{open-subscheme}{开子概形}。将\Term{sheaf}{层} $\mathcal O_{T_0}$ 看作 $U$ 上的\Term{sheaf}{层}，它便是 $\mathcal O_{U_0}$ 的商。设\Term{canonical}{典范}\Term{surjectivity}{满射}为
\[h:\mathcal O_{U_0}\longrightarrow\mathcal O_{T_0}.\]
\begin{SGAresult}{\Term{lemma}{引理}}{2.7}
\Term{canonical}{典范}映射
\[h'=\operatorname{id}_J\otimes h:J\otimes_{S_0}\mathcal O_{U_0}\longrightarrow J\otimes_{S_0}\mathcal O_{T_0}\]
\Term{decomposable}{可分解}为 $i\circ j_U$（这一\Term{resolution}{分解}显然唯一），其中 $j_U$ 是\Term{canonical}{典范}\Term{surjectivity}{满射}
\[J\otimes_{S_0}\mathcal O_{U_0}\longrightarrow J\mathcal O_U\simeq J\mathcal O_{U_0}.\]
\end{SGAresult}
须证明 $h'(K)=0$，其中 $K$ 是 $j_U$ 的核。对每个等于 $q$ 的幂的整数 $n$，有\Term{surjectivity}{满射}
\[h_n:\mathcal O_{T_0}\longrightarrow\mathcal O_{{}_nT_0},\]
而因为 ${}_nT_0$ \Term{lifting}{提升}为一个在 $S$ 上\Term{flat}{平坦}的\Term{scheme}{概形} $M_n$，\Term{canonical}{典范}\Term{morphism}{态射}
\[j_n:J\otimes_{S_0}\mathcal O_{{}_nT_0}\longrightarrow J\mathcal O_{M_n}\]
为\Term{isomorphism}{同构}。下面的\Term{commutative-diagram}{交换图}
\[
\begin{tikzcd}[column sep=large]
K\arrow[r]&J\otimes_{S_0}\mathcal O_{U_0}\arrow[r,"j_U"]\arrow[d,"h'"']&J\mathcal O_U\subset\mathcal O_U\arrow[dl,"i"]\\
&J\otimes_{S_0}\mathcal O_{T_0}\arrow[d,"h'_n"']&\\
&J\otimes_{S_0}\mathcal O_{{}_nT_0}\arrow[r,"j_n"',"\sim"]&J\mathcal O_{M_n}\subset\mathcal O_{M_n}
\end{tikzcd}
\]
证明 $h'(K)$ 对每个 $n$ 都包含在 $\operatorname{Ker}h'_n$ 中，因而包含在 $\bigcap_n\operatorname{Ker}h'_n$ 中；只需证明后者为零。但是，\Term{sheaf}{层} $\mathcal O_{T_0}$ 等于 $\mathcal O_{S_0}[\mathbb Z]$，即以 $\mathcal O_{S_0}$ 为系数的群 $\mathbb Z$ 的\Term{group-algebra}{群代数}，而 $\mathcal O_{{}_nT_0}$ 是商\Term{algebra}{代数} $\mathcal O_{S_0}[\mathbb Z/n\mathbb Z]$。
\SourcePage{263}
设 $a=\sum_{m\in\mathbb Z}a_m\otimes m$ 是 $J\otimes_{S_0}\mathcal O_{T_0}$ 的元素。因此 $a_m$ 是 $J$ 的\Term{section}{截面}，除\Term{finite}{有限}多个外均为零。取充分大的 $n$，使满足 $a_m\ne0$ 的各指标 $m$ 在 $\mathbb Z/n\mathbb Z$ 中的像两两不同。若 $a\in\operatorname{Ker}h'_n$，则必有 $a=0$。这证明 $\bigcap_n\operatorname{Ker}h'_n=0$，从而证明了 2.7。

设 $K_0$ 是 $h$ 的核；考察下面的图：
\[
\begin{tikzcd}[column sep=large]
&&&0\arrow[d]\\
&J\mathcal O_U\arrow[d,"\wr"']&&K_0\arrow[d]\\
0\arrow[r]&J\mathcal O_{U_0}\arrow[r]\arrow[d,"i"']&\mathcal O_U\arrow[r]&\mathcal O_{U_0}\arrow[r]\arrow[d,"h"]&0\\
&J\otimes_{S_0}\mathcal O_{T_0}&&\mathcal O_{T_0}\arrow[d]\\
&&&0.
\end{tikzcd}
\]
\Term{sheaf}{层} $\mathcal O_U$ 定义群 $\operatorname{Ext}^1_{\mathcal O_U}(\mathcal O_{U_0},J\mathcal O_{U_0})$ 的一个元素 $\Phi$。借助\Term{morphism}{态射} $K_0\to\mathcal O_{U_0}$ 及 $i:J\mathcal O_{U_0}\to J\otimes_{S_0}\mathcal O_{T_0}$，由 $\operatorname{Ext}^1(\ ,\ )$ 的双\Term{functor}{函子}性，从 $\Phi$ 得到 $\operatorname{Ext}^1_{\mathcal O_U}(K_0,J\otimes_{S_0}\mathcal O_{T_0})$ 中的元素 $\Psi$。由第 III 讲 4.1 及\Term{infinitesimal}{无穷小}\Term{flatness}{平坦性}判据（参见第 III 讲 4.3），存在 $U$ 的一个在 $S$ 上\Term{flat}{平坦}的\Term{subscheme}{子概形}\Term{lifting}{提升} $T_0$，当且仅当 $\Psi$ 为零。

但注意，$T_0$ \Term{affine}{仿射}，所以只需证明（第 III 讲 4.5、4.6）\Term{local-scheme}{局部}地在 $U$ 上存在一个在 $S$ 上\Term{flat}{平坦}的\Term{subscheme}{子概形}\Term{lifting}{提升} $T_0$。总之，只需证明 $\Psi$ 在\Term{sheaf}{层}
\[\mathcal E=\mathscr Ext^1_{\mathcal O_U}(K_0,J\otimes_{S_0}\mathcal O_{T_0})\]
中的像为零。仍以 $\Psi$ 表示 $\Gamma(U,\mathcal E)$ 中的这个元素。

\textbf{c) 归结到 $S$ \Term{locally-artinian}{局部 Artin}、\Term{residue-field}{剩余域}\Term{algebraically-closed}{代数闭}的情形。}
由于 $K_0$ 与 $J\otimes_{S_0}\mathcal O_{T_0}$ 都是\Term{coherent}{凝聚}\Term{sheaf}{层}，$\mathcal E$ 也是\Term{coherent}{凝聚}\Term{sheaf}{层}；此外，$\mathcal E$ 的\Term{support}{支集}包含在 $T_0$ 中，因为 $J\otimes_{S_0}\mathcal O_{T_0}$ 的\Term{support}{支集}就是如此。为证明 $\mathcal E$ 的\Term{section}{截面} $f$ 为零，只需证明对 $T_0$ 的每一点 $u$，$f$ 在 $\mathcal E$ 于 $u$ 处的\Term{fiber}{纤维}中的像为零。但\Term{coherent}{凝聚}\Term{sheaf}{层}的 $\mathscr Ext^i(\ ,\ )$ 的形成与\Term{flat}{平坦}\Term{base-extension}{基扩张}交换\footnote{编者注：参见 EGA 0$_{\mathrm{III}}$，12.3.5。}，因此归结到证明存在 $T_0\cap\operatorname{Spec}\mathcal O_u$ 的一个在 $S$ 上\Term{flat}{平坦}的\Term{lifting}{提升}。设 $s$ 是 $u$ 在 $S$ 上的\Term{projection}{投影}；于是可以将 $S$ 换成 $\operatorname{Spec}\mathcal O_s$，将 $G$ 换成 $G\times_S\operatorname{Spec}\mathcal O_s$。
\SourcePage{264}
再作一次\Term{faithfully-flat}{忠实平坦}\Term{group-extension}{扩张}，可假定 $\mathcal O_s$ 的\Term{residue-field}{剩余域}\Term{algebraically-closed}{代数闭}（EGA 0$_{\mathrm{III}}$ 10.3.1）。

设 $\mathfrak m$ 是 $\mathcal O_s$ 的\Term{maximal-ideal}{极大理想}，$S_n=\operatorname{Spec}(\mathcal O_s/\mathfrak m^n)$。假定已经证明，对每个 $n>0$，将 $T_0\times_S S_n=(T_0)_n$ \Term{lifting}{提升}为 $G_n=G\times_S S_n$ 中在 $S_n$ 上\Term{flat}{平坦}的\Term{subscheme}{子概形}的\Term{obstruction}{障碍}为零；并令 $T_n$ 为 $(T_0)_n$ 的唯一的在 $S_n$ 上\Term{flat}{平坦}且为 $G_n$ 的\Term{subtorus}{子环面}的\Term{lifting}{提升}。显然，若 $n>n'$，则
\[(T_n)\times_{S_n}S_{n'}=T_{n'}.\]
此时，若 $u$ 是 $T_0$ 中\Term{projection}{投影}到 $s$ 的一点，将下面的\Term{lemma}{引理}用于 $(T_0)_n\cap\operatorname{Spec}\mathcal O_u$ 的相容\Term{lifting}{提升}族 $(T_n\cap\operatorname{Spec}\mathcal O_u)$，便知确实存在一个在 $\mathcal O_s$ 上\Term{flat}{平坦}的 $T_0\cap\operatorname{Spec}\mathcal O_u$ 的\Term{lifting}{提升}。因此归结到证明：当 $S=S_n$ 是\Term{residue-field}{剩余域}\Term{algebraically-closed}{代数闭}的 Artin \Term{local-ring}{局部环}的谱时，$\Psi$ 为零；还须证明下面的\Term{lemma}{引理}。
\begin{SGAresult}{\Term{lemma}{引理}}{2.8}
设 $A\to B$ 是 noether \Term{local-ring}{局部环}间的\Term{local-homomorphism}{局部同态}，$\mathfrak m$ 是 $A$ 的\Term{maximal-ideal}{极大理想}，$J$ 是 $A$ 的平方为零的\Term{ideal}{理想}，$M$ 是\Term{finite-type}{有限型} $B$-模，$A_0=A/J$，$B_0=B/JB$，$M_0=M/JM$，$N_0$ 是 $M_0$ 的一个在 $A_0$ 上\Term{flat}{平坦}的 $B_0$-\Term{quotient-module}{商模}。对每个整数 $n>0$，以 $A_n,A_{0,n}$ 等表示由\Term{base-extension}{基扩张} $A\to A/\mathfrak m^n=A_n$ 得到的\Term{object}{对象}，并以 $\overline J_n$ 表示 $J$ 在 $A_n$ 中的像。对每个整数 $n>0$，设 $N_n$ 是 $M_n$ 的一个 $B_n$-\Term{quotient-module}{商模}，它在 $A_n$ 上\Term{flat}{平坦}且\Term{lifting}{提升} $N_{0,n}$，并假定当 $n\ge n'$ 时，$N_{n'}$ 由 $N_n$ 经\Term{base-extension}{基扩张} $A_n\to A_{n'}$ 得到。则存在 $M$ 的一个在 $A$ 上\Term{flat}{平坦}的 $B$-\Term{quotient-module}{商模} $N$，它\Term{lifting}{提升} $N_0$。
\end{SGAresult}
\begin{proof}[2.8 的证明]
以 $P_0$ 表示\Term{surjectivity}{满射} $M_0\to N_0$ 的核。对每个 $n>0$，有如下\Term{commutative-diagram}{交换图}，其横行均\Term{exact}{正合}：
\[
\begin{tikzcd}[column sep=small]
&&&&0\arrow[d]\\
&&&&P_{0,n}\arrow[d]\\
&0\arrow[r]&\overline J_nM_n\arrow[r]\arrow[dd]&M_n\arrow[r]\arrow[dd]&M_{0,n}\arrow[r]\arrow[dd]&0\\
\overline J_n\otimes_{A_{0,n}}M_{0,n}\arrow[urr]\arrow[drr]&&&&&\\
&0\arrow[r]&\overline J_n\otimes_{A_{0,n}}N_{0,n}\arrow[r]&N_n\arrow[r]\arrow[d]&N_{0,n}\arrow[r]\arrow[d]&0\\
&&&0&0
\end{tikzcd}\tag{$*_n$}
\]
此外，根据假设，图 $(*)_{n+1}$ 模 $\mathfrak m^n$ \Term{reduced}{约化}为 $(*)_n$。

Artin--Rees \Term{lemma}{引理}（Bourbaki，《\Term{commutative-algebra}{交换代数}》，第 3 章第 3 节推论 1）表明，由\Term{morphism}{态射}
\SourcePage{265}
\[JM\longrightarrow\overline J_nM_n\]
（相应地，$J\otimes_{A_0}M_0\to\overline J_n\otimes_{A_{0,n}}M_{0,n}$ 及 $J\otimes_{A_0}N_0\to\overline J_n\otimes_{A_{0,n}}N_{0,n}$）的核在 $JM$（相应地 $J\otimes_{A_0}M_0$ 及 $J\otimes_{A_0}N_0$）上定义的\Term{filtered}{滤过}是 $\mathfrak mB$-好\Term{filtered}{滤过}。因此，对各图 $(*)_n$ 取\Term{inverse-limit}{投射极限}，并以 $\widehat Q$ 表示 $B$-模 $Q$ 对 $\mathfrak mB$-\Term{adic-topology}{进拓扑}的\Term{separated-completion}{分离完备化}，便得到下述\Term{commutative-diagram}{交换图}，其中两个横行仍然\Term{exact}{正合}：
\[
\begin{tikzcd}[column sep=small]
&&&&0\arrow[d]\\
&&&&\widehat P_0\arrow[d]\\
&0\arrow[r]&\widehat{JM}\arrow[r]\arrow[dd]&\widehat M\arrow[r]\arrow[dd]&\widehat M_0\arrow[r]\arrow[dd]&0\\
J\widehat\otimes_{A_0}M_0\arrow[urr]\arrow[drr]&&&&&\\
&0\arrow[r]&J\widehat\otimes_{A_0}N_0\arrow[r]&\varprojlim_n N_n\arrow[r]\arrow[d]&\widehat N_0\arrow[r]\arrow[d]&0\\
&&&0&0
\end{tikzcd}\tag{$\widehat *$}
\]
但 $(B,\mathfrak mB)$ 是 Zariski 环，且 $J\otimes_{A_0}N_0$ 在 $B$ 上\Term{finite-type}{有限型}，所以它对 $\mathfrak mB$-\Term{adic-topology}{进拓扑}\Term{separated}{分离}。于是图 $\widehat{(*)}$ 表明，由\Term{surjectivity}{满射} $M_0\to N_0$ \Term{derive}{导出}的\Term{morphism}{态射}
\[J\otimes_{A_0}M_0\longrightarrow J\otimes_{A_0}N_0\]
经由 $JM$ \Term{resolution}{分解}：
\[
\begin{tikzcd}
J\otimes_{A_0}M_0\arrow[r,"\mathrm{can.}"]\arrow[dr]&JM\arrow[d]\\
&J\otimes_{A_0}N_0.
\end{tikzcd}
\]
在这些条件下，由第 III 讲 4.1 和 4.3，存在 $M$ 的一个在 $A$ 上\Term{flat}{平坦}的 $B$-\Term{quotient-module}{商模} $N$ \Term{lifting}{提升} $N_0$，当且仅当 $E=\operatorname{Ext}^1_B(P_0,J\otimes_{A_0}N_0)$ 中的某个元素 $g$ 为零。更确切地，\Term{exact-sequence}{正合列}
\[0\longrightarrow JM\longrightarrow M\longrightarrow M_0\longrightarrow0\]
定义 $F$ 中的一个元素 $f$，其中 $F$ 为 $B$-模 $\operatorname{Ext}^1_B(M_0,JM)$；而 $g$ 是 $f$ 在自然\Term{morphism}{态射} $F\to E$ 下的像，此\Term{morphism}{态射}由下述\Term{morphism}{态射}及双\Term{functor}{函子}性得到：
\[P_0\longrightarrow M_0,\qquad JM\longrightarrow J\otimes_{A_0}N_0.\]
由图 $\widehat{(*)}$ 及第 III 讲 4.1，$g$ 在 $\widehat E$ 中的像为零，其中 $\widehat E$ \Term{canonical}{典范}地等同于 $\operatorname{Ext}^1_{\widehat B}(\widehat P_0,J\widehat\otimes_{A_0}N_0)$。但 $E$ 在 $B$ 上\Term{finite-type}{有限型}，
\SourcePage{266}
所以对于 $\mathfrak mB$-\Term{adic-topology}{进拓扑}\Term{separated}{分离}，从而确有 $g=0$，这就完成了 2.8 的证明。
\end{proof}
\textbf{d) 研究 $\mathcal E$。}
现在假定 $S$ 是 Artin \Term{local-ring}{局部环} $A$ 的谱，且\Term{residue-field}{剩余域} $k$ \Term{algebraically-closed}{代数闭}。设 $A_0$ 为 $S_0$ 的环，$\mathfrak m_0$ 为 $A_0$ 的\Term{maximal-ideal}{极大理想}。由于 $S_0$ 为 Artin，$T_0$ 在 $G$ 中为闭（第 VI$_B$ 讲 1.4.2）；故可取\Term{open-subset}{开集} $U=G$，于是
\[\mathcal E=\mathscr Ext^1_{\mathcal O_G}(K_0,J\otimes_{S_0}\mathcal O_{T_0}).\]
设 $B_0$ 是\Term{affine}{仿射} $S_0$-\Term{scheme}{概形} $T_0=\mathbb G_{m,S_0}$ 的环，$\overline B_0$ 是 $T_0$ 的\Term{special-fiber}{特殊纤维} $T_0\times_{S_0}\operatorname{Spec}(k)=\mathbb G_{m,k}$ 的环。\Term{sheaf}{层} $\mathcal E$ 是一个\Term{coherent}{凝聚} $\mathcal O_{T_0}$-模，故由一个\Term{finite-type}{有限型} $B_0$-模定义，记此模为 $E$。

考察 $E$ 对 $\mathfrak m_0B_0$-进\Term{filtered-category}{滤过的}\Term{associated-graded}{相伴分次模}：
\[E_r=\mathfrak m_0^rE/\mathfrak m_0^{r+1}E.\]
因此每个 $E_r$ 都是\Term{finite-type}{有限型} $\overline B_0$-模；又由于 $S_0$ 为 Artin，当 $r$ 充分大时 $E_r=0$。

设 $g$ 是 $G$ 在 $S$ 上的一个\Term{section}{截面}，它在 $S_0$ 上诱导 $T_0$ 的\Term{section}{截面} $g_0$。由元素 $g$ 定义的\Term{translation}{平移}（为明确起见，取\Term{left-translation}{左平移}），从所考察的\Term{obstruction}{障碍}问题来看定义了“整个情形的一个\Term{automorphism}{自同构}”。特别地，$g$ 对应于\Term{sheaf}{层} $\mathcal E$ 的一个保持\Term{obstruction}{障碍} $\Psi$ \Term{invariant}{不变}的 $S_0$-\Term{automorphism}{自同构}。更确切地，$g$ 定义 $B_0$-模 $E$ 的一个\Term{semilinear}{半线性}\Term{automorphism}{自同构}（相对于群 $T_0$ 中由 $g_0$ 的\Term{translation}{平移}所定义的 $B_0$ 的 $A_0$-\Term{automorphism}{自同构}）。所以模 $\mathfrak m_0^{r+1}$ \Term{reduced}{约化}后，$g$ 定义 $E_r$ 的一个\Term{semilinear}{半线性}\Term{automorphism}{自同构}（相对于 $T_0\times_{S_0}\operatorname{Spec}(k)$ 中由 $g_0\times_{S_0}\operatorname{Spec}(k)$ 的\Term{translation}{平移}所定义的 $\overline B_0$ 的 $k$-\Term{automorphism}{自同构}）。
\begin{SGAresult}{\Term{lemma}{引理}}{2.9}
对每个整数 $r\ge0$，$E_r$ 是\Term{locally-free}{局部自由}的 $\overline B_0$-模。
\end{SGAresult}
设 $x$ 是 $T_0$ 的一点，$\kappa(x)$ 是其\Term{residue-field}{剩余域}，$(E_r)_x$ 是 $E_r$ 在 $x$ 处的“\Term{fiber}{纤维}”，即 $E_r\otimes_{\overline B_0}\kappa(x)$，$\ell(x)$ 是 $(E_r)_x$ 在 $\kappa(x)$ 上的秩，$\ell$ 是 $x$ 遍历 $T_0$ 的点时 $\ell(x)$ 的最大值。记 $L$ 为 $\operatorname{Spec}\overline B_0=\mathbb G_{m,k}$ 的最大\Term{closed-subscheme}{闭子概形}，使 $E_r$ 在其上\Term{locally-free}{局部自由}且秩为 $\ell$（TDTE IV，\Term{lemma}{引理} 3.6）。取 $L(k)$ 的一点 $\beta$（这样的点存在，因为 $L$ 在\Term{algebraically-closed-field}{代数闭域} $k$ 上\Term{finite-type}{有限型}），并取 $\mathbb G_{m,k}(k)$ 的一个阶为 $q$ 的某个幂 $n$ 的点 $\alpha$。点 $\alpha$ 在 $k$ 上有理，且根据假设，${}_nT_0$ \Term{lifting}{提升}为一个在 $S$ 上\Term{etale}{平展}的\Term{subscheme}{子概形} $M_n$，所以存在 $G$ 在 $S$ 上的\Term{section}{截面} $a$，它\Term{lifting}{提升} $\alpha$，且在 $S_0$ 上是 $T_0$ 的一个\Term{section}{截面}。上述讨论表明，$E_r$ 在 $\mathbb G_{m,k}(k)$ 的点 $\beta$ 与 $\alpha\beta$ 处的\Term{fiber}{纤维}是 $k$-\Term{isomorphism}{同构}的。但阶为 $q$ 的幂的点在 $\mathbb G_{m,k}$ 中\Term{dense}{稠密}，它们由 $\beta$ 的\Term{translation}{平移}也\Term{dense}{稠密}。由于 $L$ 在 $\mathbb G_{m,k}$ 中为闭，且 $\mathbb G_{m,k}$ \Term{reduced}{约化}，故 $L=\mathbb G_{m,k}$，从而 $E_r$ 在 $\overline B_0$ 上\Term{locally-free}{局部自由}且秩为 $\ell$。

\textbf{e) \Term{theorem}{定理} 2.1 证明的结束。}
\SourcePage{267}
以 $K_0(n)$ 表示定义\Term{closed-subscheme}{闭子概形} ${}_nT_0$ 的 $\mathcal O_{G_0}$ 的\Term{ideal-sheaf}{理想层}。\Term{sheaf}{层} $K_0$\footnote{编者注：在\Term{lemma}{引理} 2.7 之后引入。} 因而是 $K_0(n)$ 的子\Term{sheaf}{层}。为简便，置
\[R=J\otimes_{S_0}\mathcal O_{T_0},\qquad R(n)=J\otimes_{S_0}\mathcal O_{{}_nT_0}.\]
\Term{sheaf}{层} $R(n)$ 是 $R$ 的商，并有下述映射图：
\[
\begin{tikzcd}
&&&K_0\arrow[d]\\
&&&K_0(n)\arrow[d]\\
0\arrow[r]&J\mathcal O_G\arrow[r]\arrow[d]&\mathcal O_G\arrow[r]&\mathcal O_{G_0}\arrow[r]&0\\
&R\arrow[d]&&\\
&R(n)&&
\end{tikzcd}
\]
利用 $\mathscr Ext^1_{\mathcal O_G}(\ ,\ )$ 的双\Term{functor}{函子}性，得到如下\Term{commutative-diagram}{交换图}：
\[
\begin{tikzcd}[column sep=small]
\mathscr Ext^1_{\mathcal O_G}(\mathcal O_{G_0},J\mathcal O_G)\arrow[dr]&&\\
&\mathscr Ext^1_{\mathcal O_G}(K_0(n),R)\arrow[r]\arrow[d]&\mathscr Ext^1_{\mathcal O_G}(K_0,R)=\mathcal E\arrow[d]\\
&\mathscr Ext^1_{\mathcal O_G}(K_0(n),R(n))\arrow[r]&\mathscr Ext^1_{\mathcal O_G}(K_0,R(n)).
\end{tikzcd}
\]
仍以 $\Phi$ 表示由\Term{exact-sequence}{正合列}
\[0\longrightarrow J\mathcal O_G\longrightarrow\mathcal O_G\longrightarrow\mathcal O_{G_0}\longrightarrow0\]
定义的 $\mathscr Ext^1_{\mathcal O_G}(\mathcal O_{G_0},J\mathcal O_G)$ 的元素，从而 $\Psi$ 为 $\Phi$ 在 $\mathcal E$ 中的像。因为 ${}_nT_0$ \Term{lifting}{提升}为 $G$ 中一个在 $S$ 上\Term{flat}{平坦}的\Term{subscheme}{子概形} $M_n$，$\Phi$ 在\Term{sheaf}{层} $\mathscr Ext^1_{\mathcal O_G}(K_0(n),R(n))$ 中的像为零（第 III 讲 4.1）；故其在 $\mathscr Ext^1_{\mathcal O_G}(K_0,R(n))$ 中的像也为零，而这个像也是 $\Psi$ 的像。
\begin{SGAresult}{\Term{lemma}{引理}}{2.10}
\Term{canonical}{典范}\Term{morphism}{态射}
\[\mathscr Ext^1_{\mathcal O_G}(K_0,R)\otimes_{B_0}\mathcal O_{{}_nT_0}\longrightarrow\mathscr Ext^1_{\mathcal O_G}(K_0,R(n))\]
对每个整数 $n$ 都是\Term{injectivity}{单射}。
\end{SGAresult}
\SourcePage{268}
事实上，\Term{affine}{仿射}\Term{scheme}{概形} $T_0=\mathbb G_{m,S_0}$ 的环为
\[B_0=A_0[T,T^{-1}].\]
\Term{subscheme}{子概形} ${}_nT_0$ 由 $B_0$ 的下述\Term{section}{截面}的消失定义：
\[h(n)=T^n-1.\]
这个\Term{section}{截面}\Term{regular-section}{正则}（EGA 0$_{\mathrm{IV}}$ 15.2.2），并且在任意\Term{base-change}{基变换} $S'_0\to S_0$ 后仍\Term{regular-section}{正则}。因此有\Term{sheaf}{层}的\Term{exact-sequence}{正合列}
\[0\longrightarrow\mathcal O_{T_0}\xrightarrow{h(n)}\mathcal O_{T_0}\longrightarrow\mathcal O_{{}_nT_0}\longrightarrow0.\]
由于 ${}_nT_0$ 在 $S$ 上\Term{flat}{平坦}，在 $\mathcal O_{S_0}$ 上与 $J$ 张量后，得到\Term{exact-sequence}{正合列}
\[0\longrightarrow R\xrightarrow{h(n)}R\longrightarrow R(n)\longrightarrow0,\]
进而得到 $\mathscr Ext$ 的\Term{exact-sequence}{正合列}
\[\cdots\longrightarrow\mathscr Ext^1_{\mathcal O_G}(K_0,R)\xrightarrow{h(n)}\mathscr Ext^1_{\mathcal O_G}(K_0,R)\longrightarrow\mathscr Ext^1_{\mathcal O_G}(K_0,R(n)),\]
这就完成了\Term{lemma}{引理}的证明。

以上说明，对每个等于 $q$ 的幂的整数 $n$，$\Psi$ 在 $E/h(n)E$ 中的像为零。为证明 $\Psi$ 为零，只需证明若 $\Psi\in\mathfrak m_0^rE$，则 $\Psi\in\mathfrak m_0^{r+1}E$。以 $\overline\Psi$ 表示 $\Psi$ 在 $E_r$ 中的像。存在 $E$ 的一个元素 $\Psi(n)$，使
\[\Psi=\Psi(n)h(n)\qquad\text{（$n$ 等于 $q$ 的一个幂）。}\]
我们已注意到，$h(n)$ 在 $\overline B_0$ 中的像 $\overline{h(n)}$ 仍然是 $\overline B_0$-\Term{regular-section}{正则}的。由于每个 $E_s$ 都是 $\overline B_0$ 上的\Term{locally-free}{局部自由}模（2.9），$E_s$ 中乘以 $\overline{h(n)}$ 是\Term{injectivity}{单射}。由此立即\Term{pushout}{推出} $\Psi(n)\in\mathfrak m_0^rE$。以 $\overline{\Psi(n)}$ 表示其在 $E_r$ 中的像，便有关系
\[\overline\Psi=\overline{h(n)}\,\overline{\Psi(n)}\qquad\text{（$n$ 等于 $q$ 的一个幂）。}\]
这说明，$\mathbb G_{m,k}$ 上 $\overline\Psi$ 取值为零的点所成的\Term{set}{集合} $F$，包含阶为 $q$ 的幂的\Term{dense}{稠密}点集。另一方面，$F$ 为闭集（因为 $E_r$ 在 $\overline B_0$ 上\Term{locally-free}{局部自由}），且 $\mathbb G_{m,k}$ \Term{reduced}{约化}，故 $\overline\Psi$ 为零。

这就完成了\Term{theorem}{定理} 2.1 的证明。
\SGAsection{3}{用底集合刻画子环面}
\textbf{1. \Term{theorem}{定理}的陈述}

\textbf{记号。}若 $X$ 是\Term{prescheme}{预概形}，以 $\operatorname{ens}(X)$ 表示 $X$ 的\Term{underlying-set}{底集}合。若 $X,S'$ 是两个 $S$-\Term{prescheme}{预概形}，$X_{S'}=X\times_S S'$，$u:X_{S'}\to X$ 是\Term{canonical}{典范}\Term{morphism}{态射}，$E$ 是 $\operatorname{ens}(X)$ 的\Term{subset}{子集}，则以 $E_{S'}$（或 $E\times_S S'$）表示 $\operatorname{ens}(X_{S'})$ 中等于 $u^{-1}(E)$ 的\Term{subset}{子集}。
\SourcePage{269}
\begin{SGAresult}{\Term{theorem}{定理}}{3.1}
设 $S$ 是\Term{locally-noetherian}{局部 noether} \Term{prescheme}{预概形}，$G$ 是\Term{finite-type}{有限型}的 $S$-\Term{group-prescheme}{群预概形}，$q$ 是在 $S$ 上\Term{invertible}{可逆}的整数，$q>0$，$E$ 是 $\operatorname{ens}(G)$ 的一个\Term{subset}{子集}。考察关于 $E$ 的下述断言：

(i) \Term{set}{集合} $E$ 是 $G$ 的一个\Term{subtorus}{子环面} $T$ 的\Term{underlying-set}{底集}合。

(ii) a) 对 $S$ 的每一点 $s$，存在 $G_s$ 的一个\Term{subtorus}{子环面} $T_s$，使 $E\cap\operatorname{ens}(G_s)=E_s$ 是 $T_s$ 的\Term{underlying-set}{底集}合。

b) 当整数 $n$ 遍历 $q$ 的幂时，存在一族相容的（参见 2.2）$G$ 的\Term{multiplicative-type}{乘法型}\Term{group-scheme}{子群概形} $M_n$，使对 $S$ 的每一点 $s$，有
\[(M_n)_s={}_nT_s.\]

(iii) a) 同 (ii) a)。

b) \Term{set}{集合} $E$ 在 $\operatorname{ens}(G)$ 中\Term{locally-closed}{局部闭}，且 $E$ 相对于 $S$ 的\Term{fiber}{纤维}\Term{dimension}{维数}\Term{locally-constant}{局部常值}。

(iv) a) 同 (ii) a)。

b) 对每个形如\Term{residue-field}{剩余域}\Term{algebraically-closed}{代数闭}的\Term{complete-discrete-valuation-ring}{完备离散赋值环}的谱的 $S$-\Term{scheme}{概形} $S'$，$E_{S'}$ 是 $G_{S'}$ 的一个\Term{subtorus}{子环面}的\Term{underlying-set}{底集}合。

则有以下蕴涵：

A) $\mathrm{(i)}\Longleftrightarrow\mathrm{(ii)}\Longrightarrow\mathrm{(iii)}\Longrightarrow\mathrm{(iv)}$。

B) 若 $G$ 在 $S$ 上\Term{separated}{分离}，则 $\mathrm{(iii)}\Longleftrightarrow\mathrm{(iv)}$，且 $E$ 为闭集。

C) 在以下两种情形中，条件 (i)、(ii)、(iii) \Term{equivalence}{等价}（若 $G$ 在 $S$ 上\Term{separated}{分离}，则还与 (iv) \Term{equivalence}{等价}）：

第一种情形：a) \Term{prescheme}{预概形} $S$ \Term{reduced}{约化}，或者 $G$ 在 $S$ 上\Term{flat}{平坦}；且

b) 对 $S$ 的每一点 $s$，$T_s$ 是 $G_s$ 的\Term{central}{中心}\Term{torus}{环面}。

第二种情形：$S$ \Term{normal}{正规}。

此外，在以上两种情形中，以 $E$ 为底\Term{space}{空间}的\Term{torus}{环面} $T$ 唯一。
\end{SGAresult}
\textbf{注记 3.2。}a) 当 $S$ \Term{reduced}{约化}时，在 (ii) 中无须假设 $M_n$ 族相容，因为这个条件可由其他假设\Term{pushout}{推出}。事实上，若整数 $m$ 整除 $n$，\Term{subgroup}{子群} ${}_mM_n$ 和 $M_m$ 在 $S$ 上\Term{etale}{平展}，因而\Term{reduced}{约化}，且有相同的底\Term{space}{空间}，所以二者相等。

b) 若不再假定 $S$ \Term{normal}{正规}，则一般不再有 (iii)$\Rightarrow$(i)，即使 $S$ \Term{reduced}{约化}且\Term{geometrically-unibranch}{几何单分支}，$G$ 是 $S$ 上的\Term{smooth}{光滑}\Term{group-scheme}{群概形}亦然。事实上，考察 $\mathrm{SL}_{2,S}$ 中由形如
\[\begin{pmatrix}t&u\\0&t^{-1}\end{pmatrix}\]
的矩阵构成的 \Term{borel-subgroup}{Borel 子群}，其中 $S$ 是域 $k$ 上的\Term{affine}{仿射}\Term{curve}{曲线}，其环为
\[k[x,y]/(y^2-x^3).\]
按下面的方法构造\Term{set}{集合} $E$：在“$S$ 的尖点”（$x=y=0$）上取对角\Term{torus}{环面}（$u=0$）；
\SourcePage{270}
在其补\Term{open-subset}{开集}（$x\ne0$）上，取对角\Term{torus}{环面}经元素
\[\begin{pmatrix}1&y/x\\0&1\end{pmatrix}\]
\Term{conjugacy}{共轭}后得到的\Term{torus}{环面}。这样得到的\Term{set}{集合} $E$ 满足 (iii) a)；另一方面，$E$ 在 $G$ 中为闭，且以 $E$ 为\Term{underlying-set}{底集}合的\Term{reduced}{约化}\Term{subscheme}{子概形}由方程
\[\begin{cases}xu+y(t-t^{-1})=0,\\u^2-x(t-t^{-1})^2=0\end{cases}\]
定义。因此它不是 $G$ 的\Term{subtorus}{子环面}，因为尖点上的\Term{fiber}{纤维}不\Term{reduced}{约化}。

\textbf{3.1 的证明纲要。}
第一部分将证明下列“容易”的蕴涵：
\[
\begin{tikzcd}[row sep=large,column sep=large]
&\mathrm{(ii)}\arrow[dr]&\\
\mathrm{(i)}\arrow[ur]\arrow[dr]&&\mathrm{(iv)}\\
&\mathrm{(iii)}\arrow[ur]&
\end{tikzcd}
\]
以及当 $G$ 在 $S$ 上\Term{separated}{分离}时，(iv)$\Rightarrow$[(iii) 且 $E$ 为闭集]。

较困难的蕴涵分三步证明：

I) 在 C) 的第一种情形的假设下，将 (iii)$\Rightarrow$(i) 归结到 $S$ \Term{normal}{正规}的情形。

II) 当 $S$ \Term{normal}{正规}时，(iii)$\Rightarrow$(ii)。

III) (ii)$\Rightarrow$(i)。

\textbf{2. 证明 3.1 中的“容易”断言}

(i)$\Rightarrow$(ii) 及 (iii) 显然。

(iii)$\Rightarrow$(iv)。以 $S'$ 代替 $S$，可以假定 $S$ 是\Term{discrete-valuation-ring}{离散赋值环}的谱。设 $t$ 是 $S$ 的\Term{generic-point}{泛点}，$s$ 是\Term{closed-point}{闭点}。

由于 $E$ \Term{locally-closed}{局部闭}，存在 $G$ 的一个\Term{reduced}{约化}\Term{subprescheme}{子预概形}以 $E$ 为底\Term{space}{空间}，记作 $\mathsf E$。于是，$\mathsf E$ 的\Term{generic-fiber}{泛纤维} $\mathsf E_t$ 等于 $G_t$ 的秩为 $r$ 的\Term{subtorus}{子环面} $T_t$（由 (iii) a)）。设 $\mathsf E'$ 是 $\mathsf E_t$ 在 $\mathsf E$ 中的\Term{schematic-closure}{概形闭包}。\Term{prescheme}{预概形} $\mathsf E'$ \Term{irreducible}{不可约}，且其闭\Term{fiber}{纤维} $\mathsf E'_s$ \Term{nonempty}{非空}（它含有 $G_s$ 的\Term{unit-section}{单位截面}），所以 $\mathsf E'_s$ 的\Term{dimension}{维数}为 $r$（EGA IV 14.3.10）。但 $\mathsf E'_s$ 是 $\mathsf E_s$ 的\Term{closed-subscheme}{闭子概形}，后者的\Term{dimension}{维数}为 $r$（由 (iii) b)），并且\Term{irreducible}{不可约}（因为它与 $T_s$ 有相同底\Term{space}{空间}）；故 $\mathsf E'_s$ 与 $\mathsf E_s$ 的底\Term{space}{空间}相同，从而 $\mathsf E'=\mathsf E$，这证明 $\mathsf E$ 在 $S$ 上\Term{flat}{平坦}。

现设 $\mathsf E''$ 是 $\mathsf E_t$ 在 $G$ 中的\Term{schematic-closure}{概形闭包}。则 $\mathsf E''$ 是 $G$ 的一个在 $S$ 上\Term{flat}{平坦}的\Term{group-prescheme}{子群预概形}（第 VIII 讲 7.1），它包含 $\mathsf E'$，因而包含 $\mathsf E$。
\SourcePage{271}
其闭\Term{fiber}{纤维} $\mathsf E''_s$ 是 $G_s$ 的一个\Term{dimension}{维数}为 $r$ 的\Term{algebraic-subgroup}{代数子群}（同处）。由于 $T_s$ 是 $G_s$ 中的一个\Term{dimension}{维数}为 $r$ 的\Term{irreducible}{不可约}闭集，$\mathsf E_s$ 与 $\mathsf E''_s$ 的\Term{identity-component}{单位连通分支} $(\mathsf E''_s)^0$ 的\Term{underlying-set}{底集}合相同。令 $(\mathsf E'')^0$ 为 $\mathsf E''$ 的“\Term{identity-component}{单位连通分支}”，即 $\mathsf E''_s$ 的所有不包含原点的\Term{irreducible-component}{不可约分支}之并在 $\mathsf E''$ 中的补集所成的\Term{open-subgroup}{开子群}。因此有
\[E=\operatorname{ens}[(\mathsf E'')^0].\]
由于 $\mathsf E$ 和 $\mathsf E''$ \Term{reduced}{约化}，甚至有 $\mathsf E=(\mathsf E'')^0$。最终，$\mathsf E$ 是 $G$ 的一个在 $S$ 上\Term{flat}{平坦}、\Term{finite-type}{有限型}且\Term{fiber}{纤维}\Term{connected}{连通}的\Term{group-prescheme}{子群预概形}，因而\Term{separated}{分离}（第 VI$_B$ 讲 5.2）；它的\Term{generic-fiber}{泛纤维}是\Term{torus}{环面} $T_t$，其\Term{reduced}{约化}闭\Term{fiber}{纤维}是\Term{torus}{环面} $T_s$，故 $\mathsf E$ 是一个\Term{torus}{环面}（第 X 讲 8.8）。

(ii)$\Rightarrow$(iv)。仍可假定 $S$ 是\Term{discrete-valuation-ring}{离散赋值环}的谱，并沿用上述记号。$T_t$ 在 $G$ 中的\Term{schematic-closure}{概形闭包}是 $G$ 的一个在 $S$ 上\Term{flat}{平坦}的\Term{group-prescheme}{子群预概形} $\overline T$，它对每个等于 $q$ 的幂的整数 $n$ 都包含 $M_n$。因此，$\overline T$ 的闭\Term{fiber}{纤维}是 $G_s$ 的一个\Term{closed-subscheme}{闭子概形}，对每个 $n$ 都包含 $(M_n)_s$，所以包含 $T_s$。由\Term{dimension}{维数}理由，$\overline T$ 的“\Term{identity-component}{单位连通分支}”$\overline T^{\,0}$ 以 $E$ 为\Term{underlying-set}{底集}合，并像上面一样得到 $\overline T^{\,0}$ 是 $G$ 的\Term{subtorus}{子环面}。

若 $G$ 在 $S$ 上\Term{separated}{分离}，则 (iv)$\Rightarrow$[(iii) 且 $E$ 为闭集]。先证明 $E$ 为闭集。\footnote{编者注：(iv)$\Rightarrow$(iii) 这一点会在证明过程中出现。}首先证明下面的\Term{lemma}{引理}。
\begin{SGAresult}{\Term{lemma}{引理}}{3.3}
若满足 3.1 (iv) 中的条件，则 $E$ 是 $\operatorname{ens}(G)$ 的\Term{locally-constructible}{局部可构造}\Term{subset}{子集}。
\end{SGAresult}
用通常方法，归结到 $S$ 为 noether \Term{integral}{整}\Term{prescheme}{预概形}、\Term{generic-point}{泛点}为 $\eta$ 的情形。缩小 $S$ 后，可以假定（第 VI$_B$ 讲 10.10）存在 $G$ 的一个\Term{group-scheme}{子群概形} $H$，它在 $S$ 上\Term{flat}{平坦}、\Term{fiber}{纤维}\Term{connected}{连通}，且\Term{generic-fiber}{泛纤维} $H_\eta=T_\eta$。为证明 3.3，只需证明 $\operatorname{ens}(H)=E$。若 $s$ 是 $S$ 的一点，则根据 EGA II 7.1.9，存在一个 $S$-\Term{scheme}{概形} $S'$，它是\Term{discrete-valuation-ring}{离散赋值环}的谱；可假定该环\Term{complete}{完备}且\Term{residue-field}{剩余域}\Term{algebraically-closed}{代数闭}，其\Term{generic-point}{泛点} $t'$ \Term{projection}{投影}到 $\eta$，\Term{closed-point}{闭点} $s'$ \Term{projection}{投影}到 $s$。根据 (iv) b)，存在 $G_{S'}$ 的一个\Term{subtorus}{子环面} $T_{S'}$，以 $E_{S'}$ 为底\Term{space}{空间}。$G_{S'}$ 的两个\Term{group-prescheme}{子群预概形} $T_{S'}$ 和 $H_{S'}=H\times_S S'$ 都在 $S'$ 上\Term{flat}{平坦}，\Term{fiber}{纤维}\Term{connected}{连通}，且\Term{generic-fiber}{泛纤维}同为 $T_{t'}$，所以均等于 $T_{t'}$ 在 $G_{S'}$ 中的\Term{schematic-closure}{概形闭包}的\Term{identity-component}{单位连通分支}。因此 $\operatorname{ens}(H_s)=E_s$，从而 $\operatorname{ens}(H)=E$，\Term{lemma}{引理}得证。

现在已知 $E$ \Term{locally-constructible}{局部可构造}，为证明 $E$ 为闭集，只需证明 $G$ 中 $E$ 的一点的任意\Term{specialization}{特殊化}仍属于 $E$。利用通常技术，归结到 $S$ 是\Term{residue-field}{剩余域}\Term{algebraically-closed}{代数闭}的\Term{complete-discrete-valuation-ring}{完备离散赋值环}的谱的情形。但此时，$G$ 中以 $E$ 为底\Term{space}{空间}的\Term{subtorus}{子环面}（(iv) b)）在 $G$ 中为闭，因为 $G$ \Term{separated}{分离}（第 VIII 讲 7.12）。
\SourcePage{272}
\textbf{3. 3.1 证明的继续}

\textbf{I) 将 (iii)$\Rightarrow$(i)（C，第一种情形）归结到 $S$ \Term{normal}{正规的}情形。}

\textbf{a) 归结到 $S$ \Term{affine}{仿射}且\Term{reduced}{约化}的情形。}
现在假定，对 $S$ 的每一点 $s$，$E_s$ 是 $G_s$ 的一个\Term{central}{中心}\Term{subtorus}{子环面}的底\Term{space}{空间}。于是 $T$ 的唯一性由第 IX 讲 5.1 bis 得到，且 $T$ 必为\Term{central}{中心}\Term{subgroup}{子群}（同处）。由唯一性，为证明 $T$ 存在，可以假定 $S$ 为 noether \Term{affine}{仿射}\Term{prescheme}{预概形}，环为 $A$。若 $S$ 不\Term{reduced}{约化}，则根据假设，$G$ 在 $S$ 上\Term{flat}{平坦}。由 2.3，此时只需对 $S_{\mathrm{red}}$ 和 $G\times_S S_{\mathrm{red}}$ 解决问题。因此可以进一步假定 $S$ \Term{reduced}{约化}。

\textbf{b) 归结到环 $A$ 在 $\mathbb Z$ 上\Term{finite-type}{有限型}的情形。}
先证明两个\Term{lemma}{引理}。
\begin{SGAresult}{\Term{lemma}{引理}}{3.4}
设 $k$ 是域，$G$ 是 $k$-\Term{algebraic-group}{代数群}，$E$ 是 $G$ 的\Term{subset}{子集}，$k'$ 是 $k$ 的一个\Term{group-extension}{扩张}，$T_{k'}$ 是 $G_{k'}$ 的一个\Term{subtorus}{子环面}，以 $E_{k'}$ 为底\Term{space}{空间}。若 $T_{k'}$ 为\Term{central}{中心}\Term{subgroup}{子群}，或者 $k$ 是\Term{perfect-field}{完美域}，则 $E$ 是 $G_k$ 的一个\Term{subtorus}{子环面}的\Term{underlying-set}{底集}合。
\end{SGAresult}
事实上，由 fpqc \Term{descent}{下降}，只需证明 $T_{k'}$ 在 $G_{k''}$ 中的两个\Term{inverse-image}{逆像}相等，其中 $k''=k'\otimes_k k'$。但它们有相同的底\Term{space}{空间}，即 $E$ 的\Term{inverse-image}{逆像}。若 $T_{k'}$ 为\Term{central}{中心}\Term{subgroup}{子群}，则\Term{lemma}{引理}由第 IX 讲 5.1 bis 得到。若 $k$ \Term{perfect}{完美}，则 $k''$ \Term{reduced}{约化}；$T_{k'}$ 的两个\Term{inverse-image}{逆像}都在 $k''$ 上\Term{smooth}{光滑}，所以\Term{reduced}{约化}，因而相等。

\textbf{注记 3.5。}由上面的\Term{lemma}{引理}，在 3.1 (iv) 的陈述中，下列两种情形下性质 (iv) a) 可由 (iv) b) \Term{pushout}{推出}：

(1) 假定 $S$ 的各点的\Term{residue-field}{剩余域}都\Term{perfect}{完美}。

(2) 对每个如 (iv) b) 中的 $S'$，假定以 $E_{S'}$ 为底\Term{space}{空间}的\Term{torus}{环面}在 $G_{S'}$ 中为\Term{central}{中心}\Term{subgroup}{子群}。
\begin{SGAresult}{\Term{lemma}{引理}}{3.6}
设 $S$ 是\Term{affine}{仿射}\Term{scheme}{概形} $S_i$ 的\Term{inverse-limit}{投射极限}\Term{prescheme}{预概形}（参见 EGA IV 8），$H$ 是\Term{multiplicative-type}{乘法型}且\Term{finite-type}{有限型}的 $S$-\Term{group-scheme}{群概形}。则存在一个指标 $i$、一个\Term{multiplicative-type}{乘法型}且\Term{finite-type}{有限型}的 $S_i$-\Term{group-scheme}{群概形} $H_i$，以及一个 $S$-\Term{isomorphism}{同构}
\[H_i\times_{S_i}S\xrightarrow{\sim}H.\]
若 $H$ 还\Term{isotrivial}{等平凡}，则可假定 $H_i$ \Term{isotrivial}{等平凡}。
\end{SGAresult}
由于 $H$ 在 $S$ 上\Term{finite-type}{有限型}，它甚至在 $S$ 上\Term{finite-presentation}{有限呈示}（第 IX 讲 2.1 b)），所以存在一个指标 $\ell$ 及一个 $S_\ell$-\Term{group-scheme}{群概形} $H_\ell$，使 $H_\ell\times_{S_\ell}S$ \Term{isomorphism}{同构}于 $H$（第 VI$_B$ 讲第 10 节）。置 $H_i=\prod H_\ell S_iS_\ell$，于是对每个 $i\ge\ell$ 有 $H\simeq\prod H_iSS_i$。

由于 $H$ 在 $S$ 上\Term{finite-type}{有限型}，$H$ \Term{quasi-isotrivial}{拟等平凡}（第 X 讲 4.5），所以可由\Term{etale}{平展}\Term{surjectivity}{满射} $S'\to S$ \Term{trivialization}{平凡化}。利用 $S$ 的\Term{quasi-compact}{拟紧}性，容易看出存在 $S$ 的一个由\Term{finite}{有限}多个\Term{affine}{仿射}\Term{open-subset}{开集} $S_\alpha$ 构成的\Term{covering}{覆盖}，且对每个 $\alpha$，存在一个\Term{etale}{平展}、\Term{surjectivity}{满射}、\Term{finite-presentation}{有限呈示}的\Term{morphism}{态射} $S'_\alpha\to S_\alpha$，它\Term{trivialization}{平凡化} $H|_{S_\alpha}$。当 $i$ 充分大时，这个 $S$ 的\Term{covering}{覆盖} $(S_\alpha)$ 来自 $S_i$ 的一个\Term{covering}{覆盖} $(S_{i,\alpha})$（EGA IV 8）。因此，将 $S_i$ 换成 $S_{i,\alpha}$，将 $S$ 换成 $S_\alpha$，
\SourcePage{273}
便可假定 $H$ 被一个\Term{etale}{平展}、\Term{surjectivity}{满射}、\Term{finite-presentation}{有限呈示}的\Term{morphism}{态射} $S'\to S$ \Term{trivialization}{平凡化}。当 $i$ 充分大时，存在一个\Term{prescheme}{预概形} $S'_i$、一个\Term{etale}{平展}\Term{surjectivity}{满射}且\Term{finite-presentation}{有限呈示}的\Term{morphism}{态射} $S'_i\to S_i$，以及一个 $S$-\Term{isomorphism}{同构} $S'_i\times_{S_i}S\to S'$（EGA IV 17.16）。对充分大的 $j$，置
\[S'_j=S'_i\times_{S_i}S_j,\qquad H'_j=H_j\times_{S_j}S'_j,\qquad H'=H\times_S S'.\]
由 $S'$ 的选择，存在一个\Term{finite-type}{有限型}\Term{abelian-group}{阿贝尔群} $M$ 及一个 $S$-\Term{group-scheme}{群概形}\Term{isomorphism}{同构} $\mathrm D_{S'}(M)\xrightarrow{\sim}H'$。由于 $S'_i$ \Term{quasi-compact}{拟紧}且 $S'=\varprojlim_i S'_i$，由第 VI$_B$ 讲 10，存在一个指标 $j$ 及一个 $S'_j$-\Term{group-scheme}{群概形}\Term{isomorphism}{同构}
\[\mathrm D_{S'_j}(M)\xrightarrow{\sim}H'_j.\]
这就是说，$H_j$ 是\Term{quasi-isotrivial}{拟等平凡}的\Term{multiplicative-type}{乘法型}群。

当 $H$ \Term{isotrivial}{等平凡}时，采用类似的方法，使用一个\Term{finite}{有限}\Term{etale}{平展}\Term{morphism}{态射} $S'\to S$ 对 $H$ 的\Term{trivialization}{平凡化}即可。

现在可以作所说的归结 b)。$S$ 的环 $A$ 是其在 $\mathbb Z$ 上\Term{finite-type}{有限型}的子环 $A_i$ 的\Term{filtered-colimit}{滤过归纳极限}。设 $S_i=\operatorname{Spec}A_i$，$u_j:S\to S_j$ 和 $u_{jk}:S_k\to S_j$（$k\ge j$）为\Term{transition-morphism}{过渡态射}。由于 $S$ noether，且 $G$ 在 $S$ 上\Term{finite-presentation}{有限呈示}，存在一个指标 $i$ 及一个在 $S_i$ 上\Term{finite-type}{有限型}的 $S_i$-\Term{group-prescheme}{群预概形} $G_i$，使 $G$ 与 $G_i\times_{S_i}S$ 为 $S$-\Term{isomorphism}{同构}。同样，由于 $E$ 在 $G$ 中\Term{locally-closed}{局部闭}，可以假定存在 $G_i$ 的一个\Term{locally-closed-subset}{局部闭子集} $E_i$，使 $E=E_i\times_{S_i}S$（EGA IV 8.3.11）。对每个 $j>i$，记 $G_j=G_j\times_{S_i}S_j$，$E_j=E_i\times_{S_i}S_j$；令 $Q_j$ 为 $S_j$ 中使 $(E_j)_s$ 成为 $(G_j)_s$ 的一个\Term{central}{中心}\Term{subtorus}{子环面}的\Term{underlying-set}{底集}合的点 $s$ 所成的\Term{set}{集合}。

由 3.4，当 $k\ge j$ 时，$Q_k=u_{jk}^{-1}(Q_j)$；根据假设，对 $j\ge i$，有 $\operatorname{ens}(S)=u_j^{-1}(Q_j)$。此外，我断言 $Q_j$ 是\Term{ind-constructible}{归纳可构造}的（EGA IV 1.9.4）。事实上，由于 $S_j$ noether，只需证明（EGA IV 1.9.10）：若 $S$ 是\Term{generic-point}{泛点}为 $\eta$ 的 noether \Term{integral}{整}\Term{scheme}{概形}，且 $E_\eta$ 是 $G_\eta$ 的一个\Term{central}{中心}\Term{subtorus}{子环面} $T_\eta$ 的\Term{underlying-set}{底集}合，则存在 $\eta$ 的一个\Term{neighborhood}{邻域} $U$，使对 $U$ 的每一点 $s$，$E_s$ 都具有同样的性质。缩小 $S$ 后，可以假定 $G$ 的\Term{center}{中心} $Z$ \Term{representability}{可表}，并且存在 $G$ 的一个\Term{group-scheme}{子群概形} $T$，其\Term{generic-fiber}{泛纤维}为 $T_\eta$（第 VI$_B$ 讲第 10 节）。但 $\operatorname{ens}(T)$ 和 $E$ 是 $\operatorname{ens}(G)$ 的两个在\Term{generic-fiber}{泛纤维}上相同的\Term{locally-closed-subset}{局部闭子集}，所以可以找到 $\eta$ 的一个\Term{neighborhood}{邻域} $U$，使在 $U$ 上有 $\operatorname{ens}(T)=E$（EGA IV 8.3.11）。由同样的理由，可以假定在 $U$ 上，$T$ 为\Term{central}{中心}\Term{subgroup}{子群}，因为 $T$ 是\Term{torus}{环面}（3.6）。

现在已知 $Q_j$ \Term{ind-constructible}{归纳可构造}，由 EGA IV 8.3.4，当 $j$ 充分大时，$Q_j=\operatorname{ens}(S_j)$。于是可用 $S_j,G_j,E_j$ 代替 $S,G,E$，从而归结到 $S$ 是一个在 $\mathbb Z$ 上\Term{finite-type}{有限型}的\Term{reduced}{约化}\Term{affine}{仿射}\Term{scheme}{概形}的情形。

\textbf{c) 归结到 $S$ 是\Term{reduced}{约化}\Term{complete}{完备} noether \Term{local-ring}{局部环}的谱的情形。}
由于以 $E$ 为底\Term{space}{空间}的\Term{torus}{环面}唯一，通常的技术（EGA IV 8）和\Term{lemma}{引理} 3.6 允许将 $S$ 换成 $S$ 的一点的\Term{local-ring}{局部环} $A$ 的谱。设 $\widehat S$ 是 $A$ 对\Term{maximal-ideal}{极大理想}定义的\Term{topology}{拓扑}的\Term{completion}{完备化} $\widehat A$ 的谱。由于 $A$ 是 $\mathbb Z$ 上\Term{finite-type}{有限型}\Term{algebra}{代数}的\Term{localization}{局部化}，$\widehat S$ 仍然
\SourcePage{274}
\Term{reduced}{约化}（EGA IV 7.6.5）。我断言，只需在\Term{base-extension}{基扩张} $\widehat S\to S$ 后解决问题。事实上，若 $\widehat T$ 是 $G_{\widehat S}$ 中以 $E_{\widehat S}$ 为底\Term{space}{空间}的\Term{subtorus}{子环面}，则它在 $G_{\widehat S\times_S\widehat S}$ 中的两个\Term{inverse-image}{逆像}是具有相同底\Term{space}{空间}的两个\Term{central}{中心}\Term{subtorus}{子环面}，故二者相等（第 IX 讲 5.1 bis）；由 fpqc \Term{descent}{下降}，$\widehat T$ 来自 $G$ 的一个满足要求的\Term{torus}{环面} $T$（参见 3.4）。

\textbf{d) 一个\Term{descent}{下降}\Term{lemma}{引理}。}
回顾 TDTE I 中指出的\Term{finite}{有限}\Term{morphism}{态射}的下列性质。设 $S,S'$ 是两个\Term{prescheme}{预概形}，$u:S'\to S$ 是\Term{finite}{有限}\Term{morphism}{态射}。则：

(1) \Term{morphism}{态射} $u$ 是\Term{epimorphism}{满态射}，当且仅当\Term{canonical}{典范}\Term{sheaf}{层}\Term{morphism}{态射}
\[\mathcal O_S\longrightarrow u_*(\mathcal O_{S'})\]
为\Term{injectivity}{单射}。

(2) \Term{morphism}{态射} $u$ 是\Term{effective-epimorphism}{有效满态射}（第 IV 讲 1.3），当且仅当\Term{canonical}{典范}图
\[0\longrightarrow\mathcal O_S\longrightarrow u_*(\mathcal O_{S'})\rightrightarrows u_*(\mathcal O_{S'})\otimes_{\mathcal O_S}u_*(\mathcal O_{S'})\]
\Term{exact}{正合}。

(3) 若 $S$ 还 noether，且 $u$ 是\Term{epimorphism}{满态射}，则 $u$ 是\Term{finite}{有限}个\Term{finite}{有限}\Term{effective-epimorphism}{有效满态射}的复合。

现在可以证明下面的\Term{lemma}{引理}，其证明使用非\Term{flat}{平坦}\Term{descent}{下降}的技术。
\begin{SGAresult}{\Term{lemma}{引理}}{3.7}
设 $S$ 是\Term{locally-noetherian}{局部 noether} \Term{prescheme}{预概形}，$G$ 是\Term{finite-type}{有限型}的 $S$-\Term{group-prescheme}{群预概形}，$S'$ 是\Term{prescheme}{预概形}，$u:S'\to S$ 是\Term{finite}{有限}\Term{morphism}{态射}。对每个 $S$-\Term{prescheme}{预概形} $T$，以 $\mathrm M(T)$ 表示 $G_T$ 的\Term{multiplicative-type}{乘法型}\Term{subgroup}{子群}的\Term{set}{集合}（相应地，$G_T$ 的\Term{central}{中心}\Term{multiplicative-type}{乘法型}\Term{subgroup}{子群}的\Term{set}{集合}），从而 $\mathrm M$ 自然地成为定义于 $\mathrm{Sch}/S$ 上的\Term{contravariant}{反变}\Term{functor}{函子}。若 $u$ 是\Term{effective-epimorphism}{有效满态射}（相应地，\Term{epimorphism}{满态射}），则\Term{set}{集合}图
\[\mathrm M(S)\longrightarrow\mathrm M(S')\rightrightarrows\mathrm M(S'\times_S S')\tag{*}\]
\Term{exact}{正合}。
\end{SGAresult}
\begin{proof}
\textbf{i) 归结到 $u$ 是\Term{effective-epimorphism}{有效满态射}的情形。}
这里考察 $G$ 的\Term{central}{中心}\Term{multiplicative-type}{乘法型}\Term{subgroup}{子群}\Term{functor}{函子}。$\mathrm M(S)\to\mathrm M(S')$ 的\Term{injectivity}{单射}性是 $S$ 上的\Term{local-scheme}{局部}问题；承认\Term{injectivity}{单射}性后，$(*)$ 的\Term{exactness}{正合性}也成为 $S$ 上的\Term{local-scheme}{局部}问题。因此可以假定 $S$ 为 noether \Term{affine}{仿射}\Term{prescheme}{预概形}。

考察 $u:S\to S'$ 是两个\Term{finite}{有限}\Term{epimorphism}{满态射}
\[S'\xrightarrow{v}S''\xrightarrow{w}S\]
的复合的情形。我断言，若下面两个图
\[\mathrm M(S'')\longrightarrow\mathrm M(S')\rightrightarrows\mathrm M(S'\times_{S''}S')\tag{$*'$}\]
\[\mathrm M(S)\longrightarrow\mathrm M(S'')\rightrightarrows\mathrm M(S''\times_S S'')\tag{$*''$}\]
\SourcePage{275}
\Term{exact}{正合}，则 $(*)$ 也\Term{exact}{正合}。

事实上，$\mathrm M(S)\to\mathrm M(S')$ 的\Term{injectivity}{单射}性显然。若 $T'$ 是 $\operatorname{Ker}(\mathrm M(S')\rightrightarrows\mathrm M(S'\times_S S'))$ 的元素，则它当然属于 $\operatorname{Ker}(\mathrm M(S')\rightrightarrows\mathrm M(S'\times_{S''}S'))$，从而由 $(*')$ 的\Term{exactness}{正合性}，来自 $\mathrm M(S'')$ 的唯一一个元素 $T''$。只需证明 $T''$ 属于 $\operatorname{Ker}(\mathrm M(S'')\rightrightarrows\mathrm M(S''\times_S S''))$，因为由 $(*'')$ 的\Term{exactness}{正合性}，$T''$ 就可以\Term{descent}{下降}到 $S$。设 $T''_1,T''_2$ 是 $T''$ 在 $\mathrm M(S''\times_S S'')$ 中的两个\Term{inverse-image}{逆像}。由于它们是 $G_{S''\times_S S''}$ 的两个\Term{central}{中心}\Term{multiplicative-type}{乘法型}\Term{subgroup}{子群}，为证明 $T''_1=T''_2$，只需证明它们的\Term{fiber}{纤维}相同（第 IX 讲 5.1 bis）。考察\Term{commutative-diagram}{交换图}
\[
\begin{tikzcd}
S'\arrow[d,"v"']&S'\times_S S'\arrow[l,shift left]\arrow[l,shift right]\arrow[d,"v\times v"]\\
S''&S''\times_S S''\arrow[l,shift left]\arrow[l,shift right].
\end{tikzcd}
\]
\Term{morphism}{态射} $v$ 是\Term{finite}{有限}\Term{epimorphism}{满态射}，所以\Term{dominant}{支配}（上面的性质 (a)）且为\Term{closed-map}{闭映射}，因而为\Term{surjectivity}{满射}；$v\times v$ 也是如此。取 $S''\times_S S''$ 的一点 $x''$，以及其上的 $S'\times_S S'$ 的一点 $x'$。由上图的交换性，$(T''_1)_{x''}$ 与 $(T''_2)_{x''}$ 在 $G_{x'}$ 中的两个\Term{inverse-image}{逆像}相等，因此 $(T''_1)_{x''}=(T''_2)_{x''}$（EGA IV 2.2.15）。

以上讨论，再对 $u:S'\to S$ \Term{resolution}{分解}为\Term{effective-epimorphism}{有效满态射}的因子数作一个直接的归纳（上面回顾的性质 c)），表明为证明\Term{central}{中心}\Term{multiplicative-type}{乘法型}\Term{subgroup}{子群}情形中 $(*)$ 的\Term{exactness}{正合性}，可以只考虑 $u$ 为\Term{effective-epimorphism}{有效满态射}的情形。最后，再用一次第 IX 讲 5.1 bis，便知只需在 $\mathrm M$ 为 $G$ 的\Term{multiplicative-type}{乘法型}\Term{subgroup}{子群}\Term{functor}{函子}、$u$ 为\Term{effective-epimorphism}{有效满态射}时证明 3.7。

\textbf{ii) $\mathrm M(S)\to\mathrm M(S')$ 的\Term{injectivity}{单射}性。}
由于 $u:S'\to S$ 是\Term{epimorphism}{满态射}，\Term{canonical}{典范}\Term{morphism}{态射}
\[\mathcal O_S\longrightarrow u_*(\mathcal O_{S'})\]
为\Term{injectivity}{单射}；另一方面，一个\Term{multiplicative-type}{乘法型} $S$-群在 $S$ 上\Term{flat}{平坦}，因此 $\mathrm M(S)\to\mathrm M(S')$ 的\Term{injectivity}{单射}性是下面更一般\Term{lemma}{引理}的推论。
\begin{SGAresult}{\Term{lemma}{引理}}{3.8}
设 $f:X\to S$ 和 $g:S'\to S$ 是两个\Term{prescheme}{预概形}\Term{morphism}{态射}，$F$ 是\Term{quasi-coherent}{拟凝聚}的 $\mathcal O_X$-模，$X'=X\times_S S'$，$F'=F\otimes_{\mathcal O_S}\mathcal O_{S'}$。以 $Q(F)$ 表示 $F$ 的所有\Term{quasi-coherent}{拟凝聚}且在 $S$ 上\Term{flat}{平坦}的 $\mathcal O_X$-\Term{quotient-module}{商模} $\mathcal G$ 的\Term{set}{集合}，$Q(F')$ 为关于 $F',X',S'$ 的类似\Term{set}{集合}。假定 $g$ \Term{quasi-compact}{拟紧}且 $\mathcal O_S\to g_*(\mathcal O_{S'})$ 为\Term{injectivity}{单射}，则\Term{canonical}{典范}映射
\[\begin{aligned}Q(F)&\longrightarrow Q(F'),\\\mathcal G&\longmapsto\mathcal G\otimes_{\mathcal O_S}\mathcal O_{S'}\end{aligned}\]
是\Term{injectivity}{单射}。
\end{SGAresult}
事实上，可假定 $S$ \Term{affine}{仿射}，再假定 $X$ \Term{affine}{仿射}。由于\Term{morphism}{态射} $g$ \Term{quasi-compact}{拟紧}，$S'$ 是\Term{finite}{有限}多个\Term{affine}{仿射}\Term{open-subset}{开集} $S'_i$ 的并。以 $\coprod_i S'_i$ 代替 $S'$（这一操作保持 $\mathcal O_S\to g_*(\mathcal O_{S'})$ 的\Term{injectivity}{单射}性），可以假定
\SourcePage{276}
$S'$ \Term{affine}{仿射}。于是归结到下面的\Term{lemma}{引理}，其证明是直接的。
\begin{SGAresult}{\Term{lemma}{引理}}{3.9}
设 $A\to A'$ 是环的单\Term{homomorphism}{同态}，$M$ 是 $A$-模，$N=M/P$ 是 $M$ 的一个在 $A$ 上\Term{flat}{平坦}的 $A$-\Term{quotient-module}{商模}，$M'=M\otimes_A A'$，$N'=N\otimes_A A'=M'/P'$。则 $P$ 是 $P'$ 在\Term{canonical-homomorphism}{典范同态} $M\to M'$ 下的\Term{inverse-image}{逆像}（因而，知道 $N',P'$ 就知道 $N,P$）。
\end{SGAresult}
\textbf{iii) 3.7 的 $(*)$ 在 $\mathrm M(S')$ 处的\Term{exactness}{正合性}。}
设 $T'$ 是 $\operatorname{Ker}(\mathrm M(S')\rightrightarrows\mathrm M(S'\times_S S'))$ 的元素。假定已在 $S$ 为 noether \Term{local-ring}{局部环}的谱时证明 iii)。则对 $S$ 的每一点 $s$，存在 $G\times_S\operatorname{Spec}(\mathcal O_{S,s})$ 的一个\Term{multiplicative-type}{乘法型}\Term{subgroup}{子群} $T_s$，它由 $T'\times_{S'}u^{-1}\operatorname{Spec}(\mathcal O_{S,s})$ \Term{descent}{下降}而来。由第 VI$_B$ 讲第 10 节及 3.6，存在 $s$ 在 $S$ 中的一个\Term{neighborhood}{邻域} $U$，以及 $G$ 在 $U$ 上的一个\Term{multiplicative-type}{乘法型}\Term{group-scheme}{子群概形} $T$，它延拓 $T_s$。设 $U'$ 是 $U$ 在 $S'$ 中的\Term{inverse-image}{逆像}。现在得到 $G'|_{U'}$ 的两个\Term{subscheme}{子概形} $T'|_{U'}$ 和 $T\times_U U'$，它们在 $u^{-1}(\operatorname{Spec}(\mathcal O_{S,s}))$ 上相等。将 $u^{-1}(\operatorname{Spec}(\mathcal O_{S,s}))$ 看成 \Term{scheme}{概形} $u^{-1}(V)$ 的\Term{inverse-limit}{投射极限}，其中 $V$ 遍历 $s$ 在 $U$ 中的\Term{open-neighborhood}{开邻域}；由 EGA IV 8，存在 $U$ 中包含 $s$ 的\Term{open-subset}{开集} $V$，使 $T\times_U U'$ 与 $T'$ 在 $u^{-1}(V)$ 上相等。因此在所作假设下，$T'$ 可\Term{local-scheme}{局部}地在 $S$ 上\Term{descent}{下降}；但由 ii) 中证明的唯一性，$T'$ 便可整体地在 $S$ 上\Term{descent}{下降}。总之，只需在 $S$ 是 noether \Term{local-ring}{局部环}的谱时证明 iii)。

现在记 $\widehat S$ 为 $S$ 的环的\Term{completion}{完备化}的谱，令
\[S''=S'\times_S S',\qquad \widehat S'=\widehat S\times_S S',\qquad \widehat S''=\widehat S\times_S S''=\widehat S'\times_{\widehat S}\widehat S'.\]
我断言，只需证明图
\[\mathrm M(\widehat S)\longrightarrow\mathrm M(\widehat S')\rightrightarrows\mathrm M(\widehat S'')\tag{*}\]
在 $\mathrm M(\widehat S')$ 处\Term{exact}{正合}。这可由下面的\Term{commutative-diagram}{交换图}得到：按假设，第二行在 $\mathrm M(\widehat S')$ 处\Term{exact}{正合}，前两列\Term{exact}{正合}（fpqc \Term{descent}{下降}），且映射 $f$ 为\Term{injectivity}{单射}，这是将 ii) 应用于由 $S'\to S$ 经\Term{flat}{平坦}\Term{base-change}{基变换} $\widehat S\times_S\widehat S\to S$ 得到的\Term{finite}{有限}\Term{epimorphism}{满态射}
\[\widehat S'\times_{S'}\widehat S'\longrightarrow\widehat S\times_S\widehat S\]
所得的结论：
\[
\begin{tikzcd}[column sep=large]
\mathrm M(S)\arrow[r]\arrow[d]&\mathrm M(S')\arrow[r,shift left]\arrow[r,shift right]\arrow[d]&\mathrm M(S'')\arrow[d]\\
\mathrm M(\widehat S)\arrow[r]\arrow[d,shift left]\arrow[d,shift right]&\mathrm M(\widehat S')\arrow[r,shift left]\arrow[r,shift right]\arrow[d,shift left]\arrow[d,shift right]&\mathrm M(\widehat S'')\\
\mathrm M(\widehat S\times_S\widehat S)\arrow[r,"f"']&\mathrm M(\widehat S'\times_{S'}\widehat S')&
\end{tikzcd}
\]
（图追逐留给读者。）由\Term{finite}{有限}\Term{effective-epimorphism}{有效满态射}的刻画 b)，由 $S'\to S$ 经\Term{flat}{平坦}\Term{base-change}{基变换} $\widehat S\to S$ 得到的\Term{morphism}{态射} $\widehat S'\to\widehat S$ 仍为\Term{finite}{有限}\Term{effective-epimorphism}{有效满态射}。因此又归结到 $S$ 还是一个\Term{complete}{完备} noether \Term{local-ring}{局部环}的谱的情形。
\SourcePage{277}
以 $S_0$ 表示 $S$ 中以 $S$ 的\Term{closed-point}{闭点}为底\Term{space}{空间}的\Term{reduced}{约化}\Term{closed-subscheme}{闭子概形}。设 $T'$ 是 $\operatorname{Ker}(\mathrm M(S')\rightrightarrows\mathrm M(S''))$ 的元素，$T''$ 是其在 $\mathrm M(S'')$ 中的像。由于 $S'_0=S'\times_S S_0$ 在 $S_0$ 上\Term{faithfully-flat}{忠实平坦}，由 fpqc \Term{descent}{下降}，存在 $G_{S_0}$ 的一个 $S$-\Term{multiplicative-type}{乘法型}\Term{subgroup}{子群} $T_0$，使它在 $G_{S'_0}$ 中的\Term{inverse-image}{逆像}为 $T'_{S'_0}$。但 $S$ 是\Term{complete}{完备} noether \Term{local-ring}{局部环}的谱，所以存在一个\Term{multiplicative-type}{乘法型} $S$-群 $T$ 及一个 $S_0$-\Term{morphism}{态射} $u_0:T\times_S S_0\to T_0$（第 X 讲 3.3）。仍由第 X 讲 3.3，$u_0$ 在 $S'_0$ 上的\Term{inverse-image}{逆像} $u'_0$ 唯一地延拓为一个 $S'$-\Term{isomorphism}{同构} $u':T_{S'}\to T'$（注意，$S'$ 在\Term{complete}{完备}\Term{local-ring}{局部环}的谱 $S$ 上\Term{finite}{有限}，故为\Term{finite}{有限}多个\Term{complete}{完备}\Term{local-scheme}{局部}\Term{scheme}{概形}的和）。$u'$ 在 $S''$ 上的两个\Term{inverse-image}{逆像}是从 $T_{S''}$ 到 $T''$ 的两个\Term{morphism}{态射}，它们在 $S_0\times_S S''$ 上相等，所以相等（同处）。由于 $T$ 在 $S$ 上\Term{flat}{平坦}，且 $S'\to S$ 为\Term{finite}{有限}\Term{effective-epimorphism}{有效满态射}，由 TDTE I 第 8 页，图
\[\operatorname{Hom}_S(T,G)\longrightarrow\operatorname{Hom}_{S'}(T_{S'},G_{S'})\rightrightarrows\operatorname{Hom}_{S''}(T_{S''},G_{S''})\]
\Term{exact}{正合}。因此，$u'$ 来自一个\Term{prescheme}{预概形}\Term{morphism}{态射} $u:T\to G$。同样，$T\times_S T$ 在 $S$ 上\Term{flat}{平坦}，因此映射
\[\operatorname{Hom}_S(T\times_S T,G)\longrightarrow\operatorname{Hom}_{S'}(T_{S'}\times_{S'}T_{S'},G_{S'})\]
为\Term{injectivity}{单射}。立即可知 $u$ 是\Term{group-morphism}{群态射}，因为 $u'$ 是\Term{group-morphism}{群态射}。此外，$u$ 为\Term{monomorphism}{单态射}。事实上，先注意 $\operatorname{Ker}u$ 在 $S$ 上\Term{flat}{平坦}：为证明这一点，可以假定 $S$ 为 Artin \Term{local-scheme}{局部}\Term{scheme}{概形}（EGA 0$_{\mathrm{III}}$ 10.2.2），因而 $G$ \Term{separated}{分离}（第 VI$_A$ 讲 0.3）；但此时 $\operatorname{Ker}u$ 为\Term{multiplicative-type}{乘法型}（第 IX 讲 6.8），所以在 $S$ 上\Term{flat}{平坦}。由于 $\operatorname{Ker}(u)\times_S S'=\operatorname{Ker}(u')=0$，有 $\operatorname{Ker}(u)=0$（3.8）。而 $u$ 为\Term{monomorphism}{单态射}，所以是\Term{immersion}{浸入}（第 VIII 讲，注记 7.13）；其像群 $u(T)$ 确为 $\mathrm M(S)$ 的一个元素，在 $\mathrm M(S')$ 中的像为 $T'$，这就完成了 3.7 的证明。
\end{proof}
\textbf{e) I) 的证明的结束。}
由归结 c)，现在只需考虑 $S$ 为\Term{reduced}{约化}\Term{complete}{完备} noether \Term{local-ring}{局部环} $A$ 的谱的情形。设 $S'$ 为 $A$ 的\Term{normalization}{正规化} $A'$ 的谱；根据 Nagata 的\Term{theorem}{定理}（EGA 0$_{\mathrm{IV}}$ 23.1.5），$A'$ 在 $A$ 上\Term{finite}{有限}；$S'\to S$ 为\Term{epimorphism}{满态射}，因为 $A$ \Term{injectivity}{单射}到 $A'$ 中。现假定存在 $G_{S'}$ 的一个\Term{torus}{环面} $T'$，以 $E_{S'}$ 为底\Term{space}{空间}。$T'$ 在 $G_{S'\times_s S'}$ 中的两个\Term{inverse-image}{逆像}是具有相同底\Term{space}{空间}的两个\Term{central}{中心}\Term{subtorus}{子环面}，故它们相等（第 IX 讲 5 bis）；因而由 3.7，$T'$ 来自 $G_S$ 的一个\Term{central}{中心}\Term{subtorus}{子环面} $T$，它显然以 $E$ 为底\Term{space}{空间}。因此只需证明 $T'$ 存在，这将我们归结到 $S$ \Term{normal}{正规}的情形，并完成 I) 的证明。

\textbf{II) 当 $S$ \Term{normal}{正规}时证明 (iii)$\Rightarrow$(ii)。}
可以只考虑 $S$ \Term{integral}{整}的情形，设其\Term{generic-point}{泛点}为 $t$。对每个等于 $q$ 的幂的整数 $n$，以 ${}_nG$ 表示 $G$ 中的 $n$ 次幂映射的“核”\Term{subprescheme}{子预概形}。由于 $E$ 在 $G$ 中\Term{locally-closed}{局部闭}（由 (iii) b)），$\operatorname{ens}({}_nG)$ 与 $E$ 的交在 $\operatorname{ens}(G)$ 中\Term{locally-closed}{局部闭}；于是以 $E(n)$ 表示 $G$ 中以 $\operatorname{ens}({}_nG)\cap E$ 为底\Term{space}{空间}的\Term{reduced}{约化}\Term{subprescheme}{子预概形}。

下面证明\Term{structural-morphism}{结构态射} $E(n)\to S$ \Term{separated}{分离}且\Term{universally-open}{普遍开}。对这两种性质都有\Term{valuative-criterion}{赋值判据}（EGA II 7.2.3 及
\SourcePage{278}
EGA IV 14.5.8）。取一个 $S$-\Term{scheme}{概形} $S'$，它是\Term{residue-field}{剩余域}\Term{algebraically-closed}{代数闭}的\Term{complete-discrete-valuation-ring}{完备离散赋值环}的谱。我们已证明 3.1 (iii)$\Rightarrow$3.1 (iv)，所以存在 $G_{S'}$ 的一个\Term{subtorus}{子环面} $T'$，以 $E_{S'}$ 为底\Term{space}{空间}。但 ${}_nT'$ 在 $S'$ 上\Term{finite}{有限}且\Term{etale}{平展}，因而在 $S'$ 上\Term{separated}{分离}且\Term{universally-open}{普遍开}；与 ${}_nT'$ 有相同底\Term{space}{空间}的 $E(n)\times_S S'$ 也有这些性质。

另一方面，$E(n)$ 的各\Term{fiber}{纤维}有相同数目的\Term{geometric-point}{几何点}，即 $r^n$，其中 $r$ 是\Term{torus}{环面} $T_t$ 的秩。最后，\Term{generic-fiber}{泛纤维} $E(n)_t$ \Term{reduced}{约化}，所以等于 ${}_nT_t$，因而在 $\operatorname{Spec}(t)$ 上\Term{etale}{平展}。由于 $S$ \Term{normal}{正规}，由 SGA 1 10.11，$E(n)$ 是 $S$ 的一个\Term{etale}{平展}\Term{covering}{覆盖}。

若 $s$ 是 $S$ 的一点，则 $E(n)_s$ 在 $\operatorname{Spec}\kappa(s)$ 上\Term{etale}{平展}，所以\Term{reduced}{约化}，从而与\Term{multiplicative-type}{乘法型}群 ${}_nT_s$ 相等。下面证明 $E(n)$ 是 $G$ 的一个\Term{group-prescheme}{子群预概形}。设
\[m:E(n)\times_S E(n)\longrightarrow G\]
是由 $G$ 中的乘法诱导的\Term{morphism}{态射}。$m$ 的底映射 $\operatorname{ens}(m)$ 经由 $E(n)$ \Term{resolution}{分解}，故 $m$ 经由\Term{prescheme}{预概形} $E(n)$ \Term{resolution}{分解}，因为源 $E(n)\times_S E(n)$ 在 $S$ 上\Term{etale}{平展}，因而\Term{reduced}{约化}。于是由第 X 讲 4.8 a)，$E(n)$ 是一个\Term{multiplicative-type}{乘法型}\Term{subgroup}{子群}。如已指出的（3.2 a)），\Term{subgroup}{子群}族 $E(n)$ 必然相容。因此已证明当 $S$ \Term{normal}{正规}时 (iii)$\Rightarrow$(ii)。

\textbf{III) 证明 (ii)$\Rightarrow$(i)。}
事实上，我们将证明存在唯一的 $G$ 的\Term{subtorus}{子环面} $T$，以 $E$ 为底\Term{space}{空间}，并对每个等于 $q$ 的幂的整数 $n$ 满足 ${}_nT=M_n$。$T$ 的唯一性直接由第 IX 讲 4.8 b) 得到。为证明 $T$ 存在，利用唯一性，可以依次假定：

a) $S$ noether。

b) 各 $M_n$ 为\Term{central}{中心}\Term{subgroup}{子群}。事实上，只需对充分大的 $n$，以 $Z_n=\operatorname{Centr}_G(M_n)$ 代替 $G$（2.5 和 2.5 bis）。

c) $S$ 是一个\Term{local-ring}{局部环}的谱。事实上，假定对所有 $s\in S$，在\Term{base-change}{基变换} $\operatorname{Spec}(\mathcal O_{S,s})\to S$ 后问题都已解决。设 $T_s$ 是这样得到的 $G\times_S\operatorname{Spec}(\mathcal O_{S,s})$ 的\Term{subtorus}{子环面}。对每个 $s$，存在 $s$ 的一个\Term{open-neighborhood}{开邻域} $U$，及 $G|_U$ 的一个\Term{subtorus}{子环面} $T$ 延拓 $T_s$（第 VI$_B$ 讲第 10 节及 3.6）。我们已证明 3.1 (ii)$\Rightarrow$3.1 (iv)；由于 $S$ noether，$E$ 因而\Term{constructible}{可构造}（3.3）。故缩小 $U$ 后，可以假定 $\operatorname{ens}(T)=E\times_S U$（EGA IV 9.5.2）。但此时，对每个等于 $q$ 的幂的整数 $n$，${}_nT$ 与 $M_n\times_S U$ 是 $G|_U$ 的两个有相同\Term{fiber}{纤维}的\Term{multiplicative-type}{乘法型}\Term{subgroup}{子群}；由于 $M_n$ 为\Term{central}{中心}\Term{subgroup}{子群}，它们相等（第 IX 讲 5.3 bis）。总之，在所作假设下，$S$ 上\Term{local-scheme}{局部}地存在一个解，因而由唯一性存在一个整体解。

d) $S$ 是\Term{complete}{完备} noether \Term{local-ring}{局部环}的谱，且若 $s$ 是其\Term{closed-point}{闭点}，则 $T_s$ 为\Term{trivial-torus}{平凡环面}。这由 EGA 0$_{\mathrm{III}}$ 10.3.1 及 fpqc \Term{descent}{下降}得到。

e) $S$ \Term{reduced}{约化}。应用 2.2 即可。
\SourcePage{279}
f) $S$ \Term{normal}{正规}。应用 Nagata 的\Term{finiteness-theorem}{有限性定理}（EGA 0$_{\mathrm{IV}}$ 23.1.5）及 3.7 关于\Term{central}{中心}\Term{subtorus}{子环面}的情形即可。

作出这些归结后，由于 $T_s$ \Term{trivial-torus}{平凡}，$(M_n)_s$ \Term{trivial-multiplicative-group}{平凡}，所以 $M_n$ \Term{trivial-multiplicative-group}{平凡}（第 X 讲 3.3）。若 $t$ 是 $S$ 的一点，则对每个等于 $q$ 的幂的整数 $n$，$T_t$ 的\Term{subgroup}{子群} ${}_nT_t$ 都\Term{trivial-multiplicative-group}{平凡}；由第 X 讲 1.4 容易\Term{pushout}{推出} $T_t$ 本身\Term{trivial-torus}{平凡}。于是只需证明下面的\Term{lemma}{引理}。
\begin{SGAresult}{\Term{lemma}{引理}}{3.10}
在 3.1 (ii) 的假设下，再假定 $S$ noether 且\Term{normal}{正规}，并且对 $S$ 的每一点 $s$，$T_s$ 是一个\Term{trivial-torus}{平凡环面}。则存在唯一的 $G$ 的\Term{subtorus}{子环面} $T$，以 $E$ 为底\Term{space}{空间}，且对每个等于 $q$ 的幂的整数 $n$ 满足 ${}_nT=M_n$。此外，$T$ \Term{trivial-torus}{平凡}。
\end{SGAresult}
$T$ 的唯一性由 $T$ 在 $S$ 上\Term{smooth}{光滑}、因而\Term{reduced}{约化}得到。为证明存在性，可以假定 $S$ \Term{irreducible}{不可约}，\Term{generic-point}{泛点}为 $\eta$。设 $r$ 是 $T_\eta$ 的秩，$T'=\mathbb G_{m,S}^r$，$u_\eta$ 是从 $T'_\eta$ 到 $T_\eta$ 的一个\Term{isomorphism}{同构}，${}_nu_\eta$ 是 $u_\eta$ 在 ${}_nT'_\eta$ 上的\Term{restriction}{限制}。由于 ${}_nT'$ 和 $M_n$ \Term{trivial-multiplicative-group}{平凡}，存在 ${}_nu_\eta$ 的唯一延拓 ${}_nu$，它是 ${}_nT'$ 到 $M_n$ 的 $S$-\Term{isomorphism}{同构}。我断言，对 $S$ 的每一点 $s$，存在一个必然唯一的群\Term{isomorphism}{同构}
\[u_s:T'_s\xrightarrow{\sim}T_s,\]
对每个等于 $q$ 的幂的整数 $n$，它都延拓 ${}_nu_s$。事实上，取一个 $S$-\Term{scheme}{概形} $S^1$，它是\Term{residue-field}{剩余域}\Term{algebraically-closed}{代数闭}的\Term{complete-discrete-valuation-ring}{完备离散赋值环}的谱，其\Term{generic-point}{泛点} $t^1$ \Term{projection}{投影}到 $\eta$，\Term{closed-point}{闭点} $s^1$ \Term{projection}{投影}到 $s$（EGA II 7.1.9）。由 3.1 (ii)$\Rightarrow$3.1 (iv) 的证明，存在 $G_{S^1}$ 的一个\Term{subtorus}{子环面} $T^1$，使对每个等于 $q$ 的幂的整数 $n$，有 $(M_n)_{S^1}={}_nT^1$。由于 $S^1$ \Term{normal}{正规}，$T^1$ \Term{isotrivial}{等平凡}（第 X 讲 5.16）；由\Term{isotrivial}{等平凡}\Term{torus}{环面}的分类（第 X 讲 1.2）及 SGA 1 V 8.2，$T^1$ 的\Term{generic-fiber}{泛纤维}\Term{trivial-torus}{平凡}，所以它\Term{trivial-torus}{平凡}。故可将\Term{isomorphism}{同构} $u_\eta\times_\eta t^1$ 延拓为一个 $S^1$-\Term{group-scheme}{群概形}\Term{isomorphism}{同构}
\[u^1:T'\times_S S^1\xrightarrow{\sim}T^1.\]
$u^1$ 在 ${}_nT'\times_S S^1$ 上的\Term{restriction}{限制}与 ${}_nu\times_S S^1$ 按构造在\Term{generic-fiber}{泛纤维}上相等，因而相等。所以，$u^1$ 在闭\Term{fiber}{纤维}上的\Term{restriction}{限制} $u^1_{s^1}$ 在\Term{residue-field-extension}{剩余域扩张} $\kappa(s)\to\kappa(s^1)$ 后给出了所求的延拓。于是由第 IX 讲 4.8 a) 及 fpqc \Term{descent}{下降}，$u^1_{s^1}$ 在 $\kappa(s)$ 上\Term{descent}{下降}为一个满足要求的群\Term{isomorphism}{同构} $u_s:T'_s\xrightarrow{\sim}T_s$。

另一方面，$T'$ 在\Term{normal}{正规}的 $S$ 上\Term{smooth}{光滑}，所以\Term{normal}{正规}。由一个简单的\Term{rational-map}{有理映射}延拓判据（EGA IV$_4$ 20.4.6），存在一个 $S$-\Term{morphism}{态射} $u:T'\to G$，使其对 $S$ 的每一点 $s$ 在 $T'_s$ 上的\Term{restriction}{限制}为\Term{composition-morphism}{复合态射}
\[T'_s\xrightarrow{u_s}T_s\longrightarrow G_s,\]
当且仅当在每个\Term{base-change}{基变换} $S^1\to S$ 后都有这样的\Term{morphism}{态射}，其中 $S^1$ 是\Term{residue-field}{剩余域}\Term{algebraically-closed}{代数闭}的\Term{complete-discrete-valuation-ring}{完备离散赋值环}的谱。
\SourcePage{280}
在这里，一个与刚才类似的论证表明，延拓条件确实满足。以 $u$ 表示这样得到的 $S$-\Term{morphism}{态射} $T'\to G$。

下面证明 $u$ 确为\Term{group-morphism}{群态射}。设 $m_{T'}$（相应地 $m_G$）是定义 $T'$（相应地 $G$）中乘法的\Term{morphism}{态射}。须验证 $m_G\circ(u\times_S u)=u\circ m_{T'}$。这两个\Term{morphism}{态射}的相等处所构成 $T'\times_S T'$ 的一个\Term{subprescheme}{子预概形}，它包含全部\Term{fiber}{纤维}（因为 $u_S$ 是\Term{group-morphism}{群态射}），因而有与 $T'\times_S T'$ 相同的底\Term{space}{空间}，故等于 $T'\times_S T'$，因为 $T'\times_S T'$ 在 $S$ 上\Term{smooth}{光滑}，所以\Term{reduced}{约化}。

最后注意，$u$ 是\Term{monomorphism}{单态射}（因为它在各\Term{fiber}{纤维}上都是\Term{monomorphism}{单态射}），所以是\Term{immersion}{浸入}（第 VIII 讲 7.9）。于是 $T'$ 在 $u$ 下的像是 $G$ 的一个\Term{subtorus}{子环面}，以 $E$ 为\Term{underlying-set}{底集}合。

这就完成了\Term{theorem}{定理} 3.1 的证明。
